\documentclass[12pt,a4paper]{article}
\usepackage[T1]{fontenc}
\usepackage[utf8]{inputenc}
\usepackage[ngerman]{babel}
\usepackage{amsmath,amssymb,amsthm}
\usepackage{geometry}
\usepackage{titlesec}
\usepackage{fancyhdr}
\usepackage{enumitem}

\geometry{margin=2.5cm}

\title{Lehrbuch der Algebra --- Band 3, Teil 14}
\author{Heinrich Weber}
\date{}

\theoremstyle{definition}
\newtheorem{defin}{Definition}[section]
\newtheorem{satz}{Satz}[section]

\pagestyle{fancy}
\fancyhf{}
\fancyhead[C]{\small Lehrbuch der Algebra --- Band 3}
\fancyfoot[C]{\thepage}
\renewcommand{\headrulewidth}{0.4pt}

\begin{document}

\maketitle
\thispagestyle{empty}

\setcounter{section}{168}
\setcounter{page}{619}

% §169
\section{Beziehung der Teilungskörper zu dem Klassenkörper}

Wir haben in den beiden letzten Paragraphen zwei Arten von Gruppen betrachtet, von denen die erste zu den Teilungskörpern der elliptischen Funktionen, die zweite zu den Ordnungskörpern in Verbindung mit Einheitswurzeln führte. Zwischen diesen besteht eine Beziehung:

Die erste Gruppe $A$ bestand aus den Zahlen $\alpha$ in $\mathfrak{O}$, die der Bedingung
\begin{equation}
\alpha \equiv \pm 1 \pmod{\mathfrak{m}}
\end{equation}
genügten, wenn $\mathfrak{m}$ irgend ein Ideal in $\mathfrak{z}$ ist, die zweite Gruppe $A'$ bestand aus den Zahlen $\alpha$ der Ordnung $[\mathfrak{o}]$, die der Bedingung
\begin{equation}
N(\alpha) \equiv 1 \pmod{m}
\end{equation}
genügten, wenn $m$ und $\mathfrak{o}$ natürliche Zahlen bedeuten.

Nehmen wir an, es sei $\mathfrak{m}$ ein Primideal oder eine Potenz eines Primideals, so können wir $\mathfrak{o}$ und $m$ so annehmen, daß die Gruppe $A'$ in der Gruppe $A$ enthalten ist.

Denn ist $r$ eine rationale Zahl, so ist, da $\alpha$ der Ordnung $[\mathfrak{o}]$ angehört:
\begin{equation}
\alpha = r \pmod{\mathfrak{m}},
\end{equation}
und wenn $\mathfrak{o}$ durch $\mathfrak{m}$ teilbar ist:
\begin{equation}
r \equiv r^2 \pmod{\mathfrak{m}},
\end{equation}
also nach (2) und (4)
\begin{equation}
r^2 \equiv 1, \quad r(r-1) \equiv 0 \pmod{\mathfrak{m}},
\end{equation}
wir nehmen $\mathfrak{o}$ durch $\mathfrak{m}$ und $m$ durch $\mathfrak{m}$ teilbar an; dann ist nach (5), wenn $\mathfrak{m}$, wie wir angenommen haben, ein Primideal oder eine Potenz eines solchen ist, entweder
\begin{equation}
r \equiv 1 \quad \text{oder} \quad r \equiv -1 \pmod{\mathfrak{m}},
\end{equation}
und aus (4) folgt:
\begin{equation}
\alpha \equiv \pm 1 \pmod{\mathfrak{m}};
\end{equation}
folglich ist $A'$ in $A$ enthalten, und nach \S 164, 11. ist also $K(A)$ in $K(A')$ enthalten. Da nun nach \S 158 jeder Teilungskörper sich aus solchen zusammensetzen läßt, deren Teiler ein Primideal oder eine Potenz eines Primideals ist, so folgt:

\begin{satz}
Der Teilungskörper der elliptischen Funktionen ist zurückführbar auf Ordnungskörper und Kreisteilungskörper\footnote{Auf anderem Wege ist ein Teil dieses Satzes bewiesen in der Straßburger Dissertation von Daniel Bauer, \glqq Über den Teilungskörper der elliptischen Funktionen mit singulärem Modul''. Straßburg 1903.}.
\end{satz}

\newpage
\section*{FÜNFTES ALGEBRAISCHE BUCH. FUNKTIONEN.}
\addcontentsline{toc}{part}{Fünftes Buch: Algebraische Funktionen}

\setcounter{section}{169}

% §170
\section{Einleitendes}

\addcontentsline{toc}{section}{\S 170. Einleitendes}

Die Theorie der algebraischen Funktionen einer Veränderlichen ist der Ausgangspunkt der allgemeinen Untersuchungen von Abel über die neuen Transzendenten, die seitdem den Namen \glqq Abelsche Funktionen`` erhalten haben, und die höchste Verallgemeinerung der elliptischen Funktionen sind. Die Hauptprobleme dieser Theorie sind durch die Arbeiten von Riemann, Weierstraß, Clebsch, Brill und Noether zu einem gewissen Abschluß gebracht. Insbesondere hat Riemann in der Vorstellung der mehrblätterigen (Riemannschen) Flächen ein durch seine Anschaulichkeit außerordentlich wirksames Hilfsmittel für die Untersuchung dieser Funktionen geschaffen.

Alle diese Untersuchungen aber, die sich einerseits auf die Funktionentheorie, andererseits, wie bei Clebsch, Brill, Noether, auf die Methoden der rationalen Algebra (Theorie der Formen und Invarianten oder der algebraischen Kurven) stützen, müssen immer gewisse Einschränkungen machen, sie müssen gewisse \glqq Ausnahmefälle`` ausschließen und sich mit der Behandlung der sogenannten allgemeinen Fälle begnügen. Nicht unterworfen ist dieser Beschränkung die nach Analogie der Zahlentheorie von Dedekind und mir ausgearbeitete Theorie der algebraischen Funktionen, von der hier eine Übersicht gegeben werden soll\footnote{Zur Literatur über algebraische und Abelsche Funktionen sei hier erwähnt: Abel, Mém. sur une propriété générale d'une classe très-étendue de fonctions transcendantes (1826 der Pariser Akademie vorgelegt, Werke, neue Ausgabe von Sylow und Lie, Bd. I, S. 145). Riemann, Theorie der Abelschen Funktionen, Crelles Journ., Bd. 54, 1857, Werke 2. Aufl., Nr. 88. Weierstraß, Vorlesungen 1875/76, Werke, Bd. IV. Clebsch, Über die Anwendung der Abelschen Funktionen in der Geometrie, Crelles Journ., Bd. 63, 1864. Clebsch und Gordan, Theorie der Abelschen Funktionen, Leipzig 1866. Brill und Noether, Über die algebraischen Funktionen und ihre Anwendung in der Geometrie, Mathematische Annalen VII, 1874. Brill und Noether, Die Entwickelung der Theorie der algebraischen Funktionen, Bericht der Deutschen Mathematiker-Vereinigung 1894. Appell et Goursat, Théorie des fonctions algébriques, Paris 1895. Über die zahlentheoretischen Methoden: Kronecker, Über die Diskriminante algebraischer Funktionen, Crelles Journ., Bd. 91, 1881, und: Festschrift zu Kummers Doktor-Jubiläum, Crelles Journ., Bd. 92, 1881. Dedekind und Weber, Theorie der algebraischen Funktionen einer Veränderlichen, Crelles Journ., Bd. 92, 1879. Hensel und Landsberg, Theorie der algebraischen Funktionen einer Variablen, Leipzig 1902. Die Form der Theorie, die im folgenden dargestellt ist, stützt sich auf den in Bd. II, \S 153 eingeführten Begriff der Funktionale, die in dieser Form bereits in einer von Wellstein angeregten Straßburger Dissertation von Rehfeld (1906) zur Behandlung eines speziellen Problems der Funktionentheorie angewandt wurden.}.

Es würde nirgends eine Lücke bleiben, wenn wir uns dabei auf algebraische Zahlen beschränken wollten. Betrachtet man $\vartheta$, $z$ als cartesische Koordinaten in der Ebene, so stellt die Gleichung (1) eine algebraische Kurve dar, und man kann die Theorie dieser Kurven anwenden, wie es Clebsch und Gordan und Brill und Noether getan haben. Freilich haben dabei nur die reellen Werte dieser Variablen eine wirklich anschauliche Bedeutung.

\newpage
\section*{Vierundzwanzigster Abschnitt. Algebraische Funktionen einer Variablen.}
\addcontentsline{toc}{section}{Vierundzwanzigster Abschnitt. Algebraische Funktionen einer Variablen.}

% §171
\section{Definition der algebraischen Funktionen}

\addcontentsline{toc}{section}{\S 171. Definition der algebraischen Funktionen}

Eine Variable $\vartheta$ heißt eine algebraische Funktion einer unabhängigen Veränderlichen $z$, wenn die beiden Variablen durch eine algebraische Gleichung
\begin{equation}
F(\vartheta, z) = 0
\end{equation}
miteinander verbunden sind. $F$ bedeutet hierin einen Ausdruck von der Form
\begin{equation}
F(\vartheta, z) = a_0 \vartheta^n + a_1 \vartheta^{n-1} + \cdots + a_{n-1}\vartheta + a_n,
\end{equation}
dessen Koeffizienten $a_0, a_1, \ldots, a_n$ ganze rationale Funktionen von $z$ ohne gemeinschaftlichen Teiler sind. Über die konstanten Koeffizienten in diesen Ausdrücken machen wir weiter keine Voraussetzung, als daß es reelle oder komplexe Zahlen sind\footnote{Es würde nirgends eine Lücke bleiben, wenn wir uns dabei auf algebraische Zahlen beschränken wollten.}.

Wir setzen dabei die Funktion $F(\vartheta, z)$ in dem Sinne als irreduzibel voraus, daß sie nicht in Faktoren zerfallen soll, die selbst rationale Funktionen von $\vartheta$ und $z$ sind.

Jede ganze Funktion $\Phi(\vartheta, z)$ läßt sich nach Bd. I, \S 20 in irreduzible Faktoren zerlegen (abgesehen nur von konstanten Faktoren). Wir sagen, daß die Funktion $\Phi(\vartheta, z)$ dann und nur dann verschwinde, wenn unter ihren irreduziblen Faktoren die Funktion $F(\vartheta, z)$ vorkommt. Dies ist die Definition des Verschwindens, nach der, wenn man ganz korrekt sein wollte, eine jede Gleichung $\Phi(\vartheta, z) = 0$ als Kongruenz nach dem Modul $F$ aufzufassen wäre. Um nicht weitläufig zu sein, wollen wir aber diese Ausdrucksweise hier nicht brauchen.

Wenn wir die Gleichung (1) durch $a_0$ dividieren, so erhält sie die Form
\begin{equation}
f(\vartheta, z) = \vartheta^n + b_1 \vartheta^{n-1} + \cdots + b_{n-1}\vartheta + b_n = 0,
\end{equation}
worin die $b_1, b_2, \ldots, b_n$ ganze oder gebrochene rationale Funktionen von $z$ sind.

Das System aller ganzen und gebrochenen rationalen Funktionen $\Phi(\vartheta, z)$ von $\vartheta$ und $z$, in denen der Nenner nicht durch $F$ teilbar ist, also nicht verschwindet, hat die Eigenschaft, sich durch Addition, Subtraktion, Multiplikation und Division (außer durch 0) zu reproduzieren, und bildet daher einen Körper algebraischer Funktionen (Bd. I, \S 146), den ich mit $\Omega$ bezeichnen will.

Der Grad $n$ der irreduziblen Gleichung (2) oder (3), also der Gleichung niedrigsten Grades, der $\vartheta$ genügt, heißt der Grad des algebraischen Körpers in bezug auf $\vartheta$ oder auch der Grad von $\vartheta$.

Ist $\Phi(\vartheta)$ eine ganze Funktion von $\vartheta$, deren Koeffizienten ganze oder gebrochene rationale Funktionen von $z$ sind, so kann man durch Division eine Gleichung bilden:
\[\Phi(\vartheta) = f(\vartheta) \cdot q(\vartheta) + r(\vartheta),\]
worin $q(\vartheta)$ und $r(\vartheta)$ ebensolche Funktionen wie $\Phi$ sind, von denen jedoch $r(\vartheta)$ höchstens vom Grade $n-1$ ist. Wegen (3) ist dann
\[\Phi(\vartheta) = r(\vartheta).\]
Ist $\Phi(\vartheta)$ nicht durch $f(\vartheta)$ teilbar, so haben die beiden Funktionen keinen rationalen Teiler gemein, und man kann zwei Funktionen des Körpers $f_1(\vartheta)$, $g_1(\vartheta)$ so bestimmen, daß
\[f_1(\vartheta) \cdot f(\vartheta) + g_1(\vartheta) \cdot \Phi(\vartheta) = 1\]
wird (Bd. I, \S 6).

Da nun $f(\vartheta) = 0$ ist, so folgt hieraus:
\begin{equation}
g_1(\vartheta) = \frac{1}{\Phi(\vartheta)}.
\end{equation}

Man kann also jede gebrochene Funktion von $\vartheta$ durch eine ganze Funktion von $\vartheta$ ersetzen und erhält so den Satz:

\begin{satz}
Jede Funktion $\xi$ des Körpers $\Omega$ läßt sich auf eine einzige Weise in die Form setzen:
\begin{equation}
\xi = x_1 + x_2 \vartheta + x_3 \vartheta^2 + \cdots + x_n \vartheta^{n-1},
\end{equation}
deren Koeffizienten $x_1, x_2, \ldots, x_n$ rationale Funktionen von $z$ sind.
\end{satz}

Daß dies überhaupt möglich ist, schließt man aus (4) und (5); daß es nur auf eine Weise möglich ist, folgt aus der Irreduzibilität von $f(\vartheta)$; denn hätte man zwei verschiedene Darstellungen einer und derselben Größe $\xi$ in der Form (6), so würde ihre Differenz eine Gleichung für $\vartheta$ von niedrigerem als $n$tem Grade ergeben.

Wählt man $n$ Funktionen des Körpers $\Omega$ beliebig aus:
\[\eta_1, \eta_2, \ldots, \eta_n,\]
mit der einzigen Beschränkung, daß die Determinante
\[\left|\eta_i^{(k)}\right|\]
nicht identisch Null ist, so kann man durch Elimination der Potenzen von $\vartheta$ aus (6) und (7) jede Funktion $\xi$ in $\Omega$ in die Form setzen:
\begin{equation}
\xi = y_1 \eta_1 + y_2 \eta_2 + \cdots + y_n \eta_n,
\end{equation}
deren Koeffizienten $y_i$ rationale Funktionen von $z$ sind.

Ein solches System von Funktionen $\eta_1, \eta_2, \ldots, \eta_n$ heißt eine Basis des Körpers $\Omega$.

Damit ein Funktionensystem $\eta_1, \eta_2, \ldots, \eta_n$ eine Basis von $\Omega$ bilde, ist notwendig und hinreichend, daß keine Gleichung von der Form
\[y_1 \eta_1 + y_2 \eta_2 + \cdots + y_n \eta_n = 0\]
bestehe, in der die Koeffizienten $y_1, y_2, \ldots, y_n$ nicht alle identisch verschwinden.

So bilden die Funktionen
\[1, \vartheta, \vartheta^2, \ldots, \vartheta^{n-1}\]
eine Basis von $\Omega$.

\newpage
\section*{\S 172. Normen und Spuren.}
\addcontentsline{toc}{section}{\S 172. Normen und Spuren.}

Der Übergang von einer Basis $\eta_1, \eta_2, \ldots, \eta_n$ zu einer anderen $\eta'_1, \eta'_2, \ldots, \eta'_n$ geschieht durch eine lineare Substitution:
\begin{equation}
(\eta'_1, \eta'_2, \ldots, \eta'_n) = L(\eta_1, \eta_2, \ldots, \eta_n)
\end{equation}
oder kürzer geschrieben
\begin{equation}
(\eta') = L(\eta),
\end{equation}
worin $L$ die Bedeutung hat
\begin{equation}
L = \begin{pmatrix}
l_{1,1} & l_{1,2} & \cdots & l_{1,n} \\
l_{2,1} & l_{2,2} & \cdots & l_{2,n} \\
\vdots & \vdots & \ddots & \vdots \\
l_{n,1} & l_{n,2} & \cdots & l_{n,n}
\end{pmatrix}.
\end{equation}

Die Elemente $l_{i,k}$ dieser Substitution sind rationale Funktionen von $z$. Das System $(\eta')$ ist dann und nur dann wieder eine Basis, wenn die Determinante $|L|$ dieser Substitution nicht verschwindet.

\bigskip

Wählt man eine beliebige Basis $\eta_1, \eta_2, \ldots, \eta_n$, so kann man, wenn $\xi$ eine beliebige Funktion des Körpers ist, da die Produkte $\xi\eta_i$ im Körper enthalten sind, setzen:
\begin{equation}
\xi\eta_i = \gamma_{i,1}\eta_1 + \gamma_{i,2}\eta_2 + \cdots + \gamma_{i,n}\eta_n, \quad (i = 1, 2, \ldots, n)
\end{equation}
und durch Elimination der $\eta_i$ die nach Potenzen von $\xi$ geordnet, eine Gleichung
\begin{equation}
\varphi(\xi) = \xi^n + d_1\xi^{n-1} + \cdots + d_{n-1}\xi + d_n = 0
\end{equation}
erhält, in der $d_1, d_2, \ldots, d_n$ rationale Funktionen von $z$ sind. Diese Koeffizienten sind von der Wahl der Basis $\eta_i$ unabhängig. Denn die Gleichungen (1), durch die der Übergang von der Basis $\eta_i$ zu der Basis $\eta'_i$ vermittelt wird, können wir nach der Bezeichnung (9), (10) des vorigen Paragraphen so darstellen:
\begin{equation}
(\eta') = L(\eta),
\end{equation}
und demnach ergibt sich durch Übergang zu der Basis $\eta'_i$:
\[(\xi\eta') = LYL^{-1}(\eta').\]
Die Funktionen $d_1, d_2, \ldots, d_n$ sind also durch die Funktion $\xi$ vollständig bestimmt.

Die Funktion
\begin{equation}
N(\xi) = \begin{vmatrix}
\gamma_{1,1} & \gamma_{1,2} & \cdots & \gamma_{1,n} \\
\gamma_{2,1} & \gamma_{2,2} & \cdots & \gamma_{2,n} \\
\vdots & \vdots & \ddots & \vdots \\
\gamma_{n,1} & \gamma_{n,2} & \cdots & \gamma_{n,n}
\end{vmatrix} - \xi
\end{equation}
heißt die \emph{Norm} der Funktion $\xi$ und wird mit $N(\xi)$ bezeichnet. Über sie gelten folgende Sätze:

\begin{satz}
Wenn $\xi$ nicht identisch Null ist, so ist $N(\xi)$ von Null verschieden.
\end{satz}
Denn wenn die Determinante des Systems (1) verschwindet, so lassen sich rationale Funktionen $\chi_1, \chi_2, \ldots, \chi_n$, die nicht alle verschwinden, so bestimmen, daß
\[\chi_1\eta_1 + \chi_2\eta_2 + \cdots + \chi_n\eta_n = 0\]
wird, und dies fordert, da $\eta_1, \eta_2, \ldots, \eta_n$ eine Basis ist, $\xi = 0$.

\begin{satz}
Ist $\xi$ eine rationale Funktion von $z$, so ist ihre Norm die $n$te Potenz dieser Funktion.
\end{satz}
Denn ist $\xi$ rational, so reduzieren sich die Gleichungen (1) auf die Identitäten $\xi\eta_i = \xi\eta_i$. Es verschwinden also in der Determinante (5) alle Glieder mit Ausnahme der Diagonalglieder, und diese werden alle gleich $\xi$. Wir drücken diesen Satz durch die Formel aus:
\begin{equation}
N(\xi) = \xi^n.
\end{equation}

\begin{satz}
Sind $\xi, \xi'$ zwei Funktionen in $\Omega$, so gilt der Satz
\begin{equation}
N(\xi\xi') = N(\xi) \cdot N(\xi').
\end{equation}
\end{satz}
Denn ist nach (4):
\[(\xi\eta) = \Upsilon(\eta), \quad (\xi'\eta) = \Upsilon'(\eta),\]
so ist
\[(\xi\xi'\eta) = \Upsilon\Upsilon'(\eta),\]
und daraus folgt die Formel (6), weil die Determinante einer zusammengesetzten linearen Substitution, hier $\Upsilon\Upsilon'$, gleich dem Produkt der Determinanten der Komponenten ist.

\begin{satz}
Aus 2. und 3. folgt, wenn $\xi$ von Null verschieden ist:
\[N(1) = 1\]
und daraus für irgend zwei $\xi$, $\eta$:
\[N\left(\frac{\xi}{\eta}\right) = \frac{N(\xi)}{N(\eta)}.\]
\end{satz}

\begin{satz}
Ist $t$ eine unbestimmte oder variable Größe, und $\varphi(t)$ die Funktion, die für $t = \xi$ Null wird, so ist
\begin{equation}
\varphi(t) = N(t - \xi).
\end{equation}
\end{satz}
Das ergibt sich aus (1) und (2). Denn ersetzt man in (1) $\xi$ durch $t$, so ist dies gleichbedeutend damit, daß man $\gamma_{i,i}$ durch $\gamma_{i,i} - t$ ersetzt und die übrigen $\gamma_{i,k}$ ungeändert läßt.

Zerlegen wir die Funktion $\varphi(t)$, die eine rationale Funktion von $t$ und $z$ ist, in ein Produkt von irreduziblen Faktoren $\varphi_1(t), \varphi_2(t), \ldots$, so muß einer dieser Faktoren für $t = \xi$ verschwinden. Sei dies $\varphi_1(t)$, und
\begin{equation}
\varphi_1(t) = t^e + c_1 t^{e-1} + \cdots + c_e
\end{equation}
sei vom Grade $e$.

Es ist dann $\varphi_1(\xi) = 0$ die Gleichung niedrigsten Grades, der $\xi$ genügt, und aus $\xi$ leitet man einen Körper $\Omega_1$ ab, der in $\Omega$ enthalten ist und den Grad $e$ hat. Jede Zahl $\eta$ des Körpers $\Omega_1$ kann, und zwar auf eine Art, in der Form dargestellt werden:
\begin{equation}
\eta = \eta_1 + \eta_2\xi + \cdots + \eta_e\xi^{e-1}.
\end{equation}

Sei ferner
\begin{equation}
f(\vartheta) = 0
\end{equation}
die Gleichung niedrigsten Grades mit Koeffizienten in $\Omega_1$, der $\vartheta$ genügt, also $f$ der Relativgrad von $\Omega$ in bezug auf $\Omega_1$.

Durch die $ef$ Funktionen
\begin{equation}
\xi^i \vartheta^k \quad (i = 0, 1, \ldots, e-1; \; k = 0, 1, \ldots, f-1)
\end{equation}
kann jede Funktion des Körpers linear ausgedrückt werden, und zwischen ihnen besteht keine lineare Gleichung mit rational von $z$ abhängigen Koeffizienten. Denn sonst würde $\vartheta$ oder $\xi$ aus einer Gleichung von niedrigerem Grade als $f$ oder $e$ entstehen. Demnach bilden die Funktionen (12) eine Basis des Körpers $\Omega$, und es folgt
\begin{equation}
n = ef.
\end{equation}

Es sind also $e$ und $f$ Teiler von $n$.

Bedeutet
\begin{equation}
\varepsilon_1, \varepsilon_2, \ldots, \varepsilon_e
\end{equation}
eine Basis des Körpers $\Omega_1$, so bilden die $n$ Größen:
\begin{equation}
\varepsilon_i \vartheta^k \quad (i = 1, \ldots, e; \; k = 0, \ldots, f-1)
\end{equation}
eine Basis von $\Omega$.

Bildet man die Substitution der $e$ Größen:
\[\xi\varepsilon_i = \sum_j \gamma_{i,j}\varepsilon_j,\]
so ist
\begin{equation}
N_1(\xi) = |\gamma_{i,j}|
\end{equation}
die Norm von $\xi$ im Körper $\Omega_1$, und wenn man die Substitution für $(\xi \cdot \vartheta^k)$ im Körper $\Omega$ mit der Basis (15) bildet, so ergibt sich
\[\Upsilon_i(\xi \cdot \vartheta^k) = \Upsilon_1(\xi) \cdot \Upsilon(\vartheta^k),\]
woraus
\begin{equation}
N = N_1^f.
\end{equation}

Aus der Substitutionsdeterminante $|\Upsilon_i|$ leitet man die Funktion $\varphi_1(t)$ nach der Formel (2) ab, und es folgt also nach (18):
\begin{equation}
N(t - \xi) = \varphi_1(t)^f,
\end{equation}
und daraus:

\begin{satz}
Die Funktion $\varphi(t) = N(t - \xi)$ ist entweder irreduzibel oder eine Potenz einer irreduziblen Funktion.
\end{satz}

Ist der Exponent $f$ dieser Potenz größer als 1, so ist der Körper $\Omega$ imprimitiv, $\Omega_1$ ist ein Teilkörper, und $\xi$ ist ein Körper vom Grade $f$ über $\Omega_1$ (Bd. I, \S 151). Ist $f = 1$, $e = n$, so heißt $\xi$ eine primitive Funktion des Körpers $\Omega$, und $\Omega_1$ ist mit $\Omega$ identisch.

Wir kehren zu der Gleichung (2) oder (3) zurück, die, wie wir gesehen haben, von der Wahl der Basis unabhängig ist, und definieren den Koeffizienten
\begin{equation}
S(\xi) = -d_1 = \gamma_{1,1} + \gamma_{2,2} + \cdots + \gamma_{n,n}
\end{equation}
als die \emph{Spur} der Funktion $\xi$. Für die Spur ergeben sich aus der Definition die folgenden fundamentalen Sätze, in denen $\xi$, $\eta$ irgend zwei Funktionen in $\Omega$, und $x$ eine rationale Funktion bedeutet:
\begin{align}
S(\xi + \eta) &= S(\xi) + S(\eta), \\
S(x\xi) &= xS(\xi), \\
S(x) &= nx, \\
S(\xi\eta) &= S(\eta\xi).
\end{align}
Alle Spuren sind rationale Funktionen von $z$.

\newpage
\section*{\S 173. Diskriminanten.}
\addcontentsline{toc}{section}{\S 173. Diskriminanten.}

Ist $(\eta_1, \eta_2, \ldots, \eta_n)$ ein System von $n$ Funktionen im Körper $\Omega$, so sind die $n^2$ Spuren $S(\eta_i\eta_k)$ rationale Funktionen von $z$.

Wir definieren als \emph{Diskriminante} des Systems $(\eta_i)$ die Determinante
\[\Delta(\eta_i) = \begin{vmatrix}
S(\eta_1\eta_1) & S(\eta_1\eta_2) & \cdots & S(\eta_1\eta_n) \\
S(\eta_2\eta_1) & S(\eta_2\eta_2) & \cdots & S(\eta_2\eta_n) \\
\vdots & \vdots & \ddots & \vdots \\
S(\eta_n\eta_1) & S(\eta_n\eta_2) & \cdots & S(\eta_n\eta_n)
\end{vmatrix},\]
wofür wir auch kürzer $\Delta(\eta_i)$ schreiben.

Wir beweisen den Satz:

\begin{satz}
Die Diskriminante $\Delta(\eta_i)$ ist dann und nur dann von Null verschieden, wenn $(\eta_i)$ eine Basis des Körpers $\Omega$ ist.
\end{satz}

Nehmen wir, um ihn zu beweisen, zunächst an, daß $\Delta(\eta_i) = 0$ sei. Dann kann man nach einem elementaren Determinantensatz die rationalen Funktionen $\chi_1, \chi_2, \ldots, \chi_n$, ohne daß sie alle Null sind, so bestimmen, daß
\begin{equation}
\chi_1 S(\eta_1\eta_k) + \chi_2 S(\eta_2\eta_k) + \cdots + \chi_n S(\eta_n\eta_k) = 0 \quad (k = 1, 2, \ldots, n)
\end{equation}
ist. Bedeutet $\xi$ ein beliebiges System rationaler Funktionen, und setzt man
\begin{equation}
\xi = \chi_1\eta_1 + \chi_2\eta_2 + \cdots + \chi_n\eta_n,
\end{equation}
so folgt aus (2) durch Multiplikation mit $\chi_k$ und Summation:
\begin{equation}
S(\xi\eta_k) = 0.
\end{equation}
Ist nun $(\eta_i)$ eine Basis, so kann $\xi$ jede beliebige Funktion in $\Omega$, also auch $1/\xi$ sein, und dann gibt die Formel (4) das widersprechende Resultat $S(1) = 0$; also kann die Diskriminante einer Basis $(\eta_i)$ nicht verschwinden.

Um auch das Umgekehrte zu beweisen, nehmen wir an, es sei $(\eta_i)$ eine Basis und
\begin{equation}
(\eta'_i) = X(\eta_i)
\end{equation}
eine lineare Substitution. Das System $(\eta'_i)$ ist dann und nur dann gleichfalls eine Basis, wenn die Determinante $|X|$ dieser Substitution von Null verschieden ist (\S 171). Setzt man nach (5)
\[\eta'_i = \sum_j x_{i,j}\eta_j,\]
so wird
\[S(\eta'_i\eta'_k) = \sum_{r,s} x_{i,r}x_{k,s}S(\eta_r\eta_s),\]
und indem man daraus die Determinante bildet:
\begin{equation}
\Delta(\eta'_1, \ldots, \eta'_n) = |X|^2 \Delta(\eta_1, \ldots, \eta_n).
\end{equation}
Ist also $\Delta(\eta_i)$ von Null verschieden, so kann $|X|$ nicht verschwinden und $(\eta'_i)$ ist eine Basis. Damit ist 7. bewiesen.

Aus (6) ergibt sich noch nach der Definition der Norm in \S 172, wenn man $\eta_i = \xi\eta_i$ setzt, und $\xi$ eine beliebige Funktion in $\Omega$ ist:
\begin{equation}
\Delta(\xi\eta_i) = N(\xi)^2 \Delta(\eta_i).
\end{equation}

\newpage
\section*{\S 174. Die Potenzsummen.}
\addcontentsline{toc}{section}{\S 174. Die Potenzsummen.}

Die Spuren der Potenzen von $\vartheta$
\begin{equation}
s_r = S(\vartheta^r)
\end{equation}
sind nichts anderes als die Potenzsummen, die nach den Newtonschen Formeln berechnet werden können. Da wir aber hier nicht von den \glqq Wurzeln`` der Gleichung $f(\vartheta) = 0$ sprechen dürfen, die wir noch nicht haben, so müssen diese Formeln direkt aus der Definition der Spur abgeleitet werden.

Ist $f(\vartheta) = 0$ die den Körper $\Omega$ definierende Gleichung, so setzen wir
\begin{equation}
f(t) = t^n + a_1t^{n-1} + \cdots + a_{n-1}t + a_n
\end{equation}
und bilden den Quotienten:
\[\frac{f(t)}{t - \vartheta} = \eta_0 t^{n-1} + \eta_1 t^{n-2} + \cdots + \eta_{n-2}t + \eta_{n-1},\]
worin
\begin{equation}
\begin{aligned}
\eta_0 &= 1, \\
\eta_1 &= \vartheta + a_1, \\
\eta_2 &= \vartheta^2 + a_1\vartheta + a_2, \\
&\vdots \\
\eta_{n-1} &= \vartheta^{n-1} + a_1\vartheta^{n-2} + \cdots + a_{n-1}.
\end{aligned}
\end{equation}

Zwischen diesen Funktionen bestehen die Relationen:
\begin{equation}
\begin{aligned}
\eta_0 &= 1, \\
\eta_1 &= \vartheta\eta_0 + a_1, \\
\eta_2 &= \vartheta\eta_1 + a_2, \\
&\vdots \\
\eta_{n-1} &= \vartheta\eta_{n-2} + a_{n-1}, \\
0 &= \vartheta\eta_{n-1} + a_n,
\end{aligned}
\end{equation}
und die $\eta_0, \eta_1, \ldots, \eta_{n-1}$ bilden eine Basis von $\Omega$, da man durch (4) die Potenzen von $\vartheta$ linear durch die $\eta_i$ ausdrücken kann.

Ist also $\xi$ eine beliebige Funktion in $\Omega$, so kann man setzen:
\begin{equation}
\xi = \psi_0\eta_0 + \psi_1\eta_1 + \cdots + \psi_{n-1}\eta_{n-1}.
\end{equation}

Wir leiten aus $\psi_0, \psi_1, \ldots, \psi_{n-1}$ eine Fortsetzung der Reihe dieser rationalen Funktionen $\psi_n, \psi_{n+1}, \psi_{n+2}, \ldots$ ab, indem wir setzen:
\begin{equation}
a_n\psi_r + a_{n-1}\psi_{r+1} + \cdots + a_1\psi_{r+n-1} + \psi_{r+n} = 0.
\end{equation}

Dann ergibt sich aus (4) und (5):
\[\xi\eta_0 = \psi_0\eta_0 + \psi_1\eta_1 + \cdots + \psi_{n-1}\eta_{n-1},\]
\[\xi\eta_1 = \psi_0(\vartheta\eta_0 + a_1) + \psi_1(\vartheta\eta_1 + a_2) + \cdots = \psi_1\eta_0 + \psi_2\eta_1 + \cdots + \psi_n\eta_{n-1},\]
und allgemein für ein positives $r$:
\begin{equation}
\xi\eta_r = \psi_r\eta_0 + \psi_{r+1}\eta_1 + \cdots + \psi_{r+n-1}\eta_{n-1},
\end{equation}
wie man aus (4) und (6) durch den Schluß von $r$ auf $r+1$ leicht bestätigt. Ordnet man diese Summe nach Potenzen von $\vartheta$, so ergibt sich ein Ausdruck von der Form
\[\xi\eta_r = \chi_{0,r}\eta_0 + \chi_{1,r}\eta_1 + \cdots + \chi_{n-1,r}\eta_{n-1},\]
worin
\begin{align*}
\chi_{0,r} &= \psi_r a_{n-1} + \psi_{r+1}a_{n-2} + \psi_{r+2}a_{n-3} + \cdots + \psi_{r+n-1}, \\
\chi_{1,r} &= \psi_r a_{n-2} + \psi_{r+1}a_{n-3} + \cdots + \psi_{r+n-2}, \\
&\vdots \\
\chi_{n-2,r} &= \psi_r a_1 + \psi_{r+1}, \\
\chi_{n-1,r} &= \psi_r.
\end{align*}

Nach der Definition der Spur [\S 172, (20)] ist aber
\[S(\xi\eta_r) = \chi_{r,r} = \psi_r a_{n-1} + \psi_{r+1}a_{n-2} + \cdots + \psi_{r+n-1},\]
und wenn man $r = 0, 1, \ldots, n-1$ setzt und addiert:
\begin{equation}
S(\xi) = n\psi_{n-1} + (n-1)\psi_{n-2}a_1 + (n-2)\psi_{n-3}a_2 + \cdots + \psi_0a_{n-1}.
\end{equation}

Will man diesen Ausdruck auf $\xi = \vartheta^k$ anwenden, so hat man $\psi_0 = 1$, die übrigen $\psi = 0$ zu setzen und erhält
\[s_k = S(\vartheta^k) = -ka_k + \cdots\]
und folglich, solange $r < n$ ist, nach (4) mit der Bezeichnung (1):
\begin{equation}
s_r + a_1s_{r-1} + a_2s_{r-2} + \cdots + a_{r-1}s_1 + ra_r = 0,
\end{equation}
und wenn man die Spur von $\vartheta^r f(\vartheta) = 0$ nimmt:
\begin{equation}
a_ns_r + a_{n-1}s_{r+1} + \cdots + a_1s_{r+n-1} + s_{r+n} = 0.
\end{equation}

Dies sind die Newtonschen Formeln (Bd. I, \S 46).

Bezeichnen wir mit
\[f'(t) = nt^{n-1} + (n-1)a_1t^{n-2} + \cdots + 2a_{n-2}t + a_{n-1}\]
die abgeleitete Funktion von $f(t)$, multiplizieren die Gleichung (12) mit $\vartheta^r t^{n-r-1}$ und summieren von $r = 0$ bis $r = n-1$, so folgt:
\begin{equation}
f'(\vartheta) = n\eta_0s_{n-1} + (n-1)\eta_1s_{n-2} + \cdots + 2\eta_{n-2}s_1 + \eta_{n-1}s_0,
\end{equation}
und wenn man für irgend ein positives $k$ die Gleichungen (12) mit $\vartheta^{r+k+n-1}$ multipliziert und noch so viele von den Gleichungen (13) hinzunimmt, bis $n - r + k - 1$ anfängt, negativ zu werden, so ergibt sich:
\begin{equation}
\vartheta^k f'(\vartheta) = \eta_0s_{n+k-1} + \eta_1s_{n+k-2} + \cdots + \eta_{n-1}s_k.
\end{equation}

Um die Norm von $f'(\vartheta)$ zu bilden, hätte man in diesen Gleichungen zunächst die Basis $\eta_0, \eta_1, \ldots, \eta_{n-1}$ durch die Substitution (2) durch $(1, \vartheta, \vartheta^2, \ldots, \vartheta^{n-1})$ auszudrücken, dann die Gleichung (15) für $k = 0, 1, \ldots, n-1$ zu bilden und die Determinante dieses Systems zu nehmen. Statt dessen kann man die Determinante in bezug auf die $n$ nehmen, und dann mit der Substitutionsdeterminante
\begin{equation}
\begin{vmatrix}
a_{n-1} & a_{n-2} & \cdots & a_1 & 1 \\
a_{n-2} & a_{n-3} & \cdots & 1 & 0 \\
\vdots & \vdots & \ddots & \vdots & \vdots \\
a_1 & 1 & \cdots & 0 & 0 \\
1 & 0 & \cdots & 0 & 0
\end{vmatrix}
= (-1)^{\frac{n(n-1)}{2}}
\end{equation}
multiplizieren. Dadurch erhält man:
\[Nf'(\vartheta) = \begin{vmatrix}
s_0 & s_1 & \cdots & s_{n-1} \\
s_1 & s_2 & \cdots & s_n \\
\vdots & \vdots & \ddots & \vdots \\
s_{n-1} & s_n & \cdots & s_{2n-2}
\end{vmatrix}.\]

Diese Determinante ist aber nach der Definition der Diskriminante (\S 173):
\[= (-1)^{\frac{n(n-1)}{2}} \Delta(1, \vartheta, \vartheta^2, \ldots, \vartheta^{n-1}),\]
und wir haben also:
\begin{equation}
Nf'(\vartheta) = (-1)^{\frac{n(n+1)}{2}} \Delta(1, \vartheta, \vartheta^2, \ldots, \vartheta^{n-1}).
\end{equation}

In der Betrachtung dieses Abschnittes haben wir, dem Hauptziel der Untersuchung entsprechend, die Koeffizienten $a_0, a_1, \ldots, a_n$ als rationale Funktionen von $z$ betrachtet. Alles bleibt aber vollständig ungeändert, wenn wir für diese Größen irgend einen Rationalitätsbereich festsetzen.

\newpage
\section*{\S 175. Ganze Funktionen von $z$.}
\addcontentsline{toc}{section}{\S 175. Ganze Funktionen von $z$.}

Jede Funktion $\xi$ des Körpers $\Omega$ genügt, wie wir gesehen haben, einer Gleichung niedrigsten Grades:
\begin{equation}
\xi^n + b_1\xi^{n-1} + \cdots + b_{n-1}\xi + b_n = 0,
\end{equation}
deren Koeffizienten $b_1, b_2, \ldots, b_n$ rationale Funktionen von $z$ sind.

\begin{defin}
Wenn diese Koeffizienten ganze rationale Funktionen von $z$ sind, so heißt $\xi$ eine \emph{ganze algebraische} oder kurz eine \emph{ganze Funktion} von $z$.
\end{defin}

Über die ganzen algebraischen Funktionen gelten die nämlichen Sätze, wie über die ganzen Zahlen (Bd. II, \S 149).

Setzen wir
\begin{equation}
g(t) = N(t - \xi),
\end{equation}
so ist, wie wir in \S 172 gesehen haben,
\begin{equation}
N(t - \xi) = \varphi(t)^f
\end{equation}
eine ganze Potenz von $\varphi(t)$. Wenn wir daher $N(t - \xi)$ nach Potenzen von $t$ ordnen, so werden die Koeffizienten alle wieder ganze rationale Funktionen, also insbesondere:

\begin{satz}
Die Norm und die Spur einer ganzen Funktion $\xi$ sind ganze rationale Funktionen von $z$.
\end{satz}

\begin{satz}
Eine rationale Funktion von $z$ ist nur dann eine ganze algebraische Funktion, wenn sie eine ganze rationale Funktion ist.
\end{satz}
Denn ist $\xi = -b$ rational, so ist $\xi + b = 1$ und $\xi + b = 0$ die Gleichung niedrigsten Grades für $\xi$, also $\xi$ nur dann ganz, wenn $b$ ganz und rational ist.

\begin{satz}
Jede Funktion $\eta$ in $\Omega$ kann durch Multiplikation mit einer von Null verschiedenen rationalen Funktion von $z$ in eine ganze algebraische Funktion verwandelt werden.
\end{satz}
Denn jede Funktion $\eta$ genügt einer Gleichung niedrigsten Grades:
\begin{equation}
b_0\eta^n + b_1\eta^{n-1} + \cdots + b_{n-1}\eta + b_n = 0,
\end{equation}
und wenn man mit $b_0^{n-1}$ multipliziert und $\xi = b_0\eta$ setzt, so erhält man eine Gleichung von der Form (1).

\begin{satz}
Eine Funktion $\xi$ des Körpers $\Omega$ ist eine ganze Funktion, wenn sie irgend einer Gleichung
\begin{equation}
\Psi(\xi) = \xi^m + c_1\xi^{m-1} + \cdots + c_{m-1}\xi + c_m = 0
\end{equation}
genügt, in der $c_1, c_2, \ldots, c_m$ ganze rationale Funktionen sind, wenn dies auch nicht die Gleichung niedrigsten Grades ist.
\end{satz}

Denn ist $\varphi(\xi) = 0$ die Gleichung niedrigsten Grades für $\xi$, so muß $\Psi(t)$ für ein variables $t$ durch $\varphi(t)$ teilbar sein. Das ist einfach der Inhalt der Gleichung (5). Es ist also
\[\Psi(t) = \varphi(t) \cdot \chi(t),\]
worin $\chi(t)$ ebenfalls eine ganze Funktion von $t$ ist. Wenn nun, wie vorausgesetzt, $\Psi(t)$ ganze Koeffizienten hat, so müssen nach einem Satz von Gauss auch die Koeffizienten von $\varphi(t)$ und von $\chi(t)$ ganze Funktionen sein. Vgl. Bd. I, \S 2, \S 20, \S 148\footnote{Bei diesem Beweis braucht nicht die Zerlegbarkeit der ganzen rationalen Funktionen in lineare Faktoren vorausgesetzt zu werden, sondern nur die Zerlegbarkeit in irreduzible Faktoren. Man kann also für die Zahlenkoeffizienten noch einen beliebigen Rationalitätsbereich festsetzen. Wir werden aber in der Folge die Zerlegbarkeit einer ganzen rationalen Funktion von $z$ in lineare Faktoren, d.\,h. den Fundamentalsatz der Algebra, doch voraussetzen müssen.}.

Hieraus ergibt sich der folgende Hauptsatz:

\begin{satz}
Summe, Differenz, Produkt zweier ganzer Funktionen in $\Omega$ sind wieder ganze Funktionen.
\end{satz}

Sind nämlich $\xi'$, $\xi''$ zwei ganze Funktionen in $\Omega$, die den Gleichungen genügen:
\begin{equation}
\xi'^n + b_1\xi'^{n-1} + \cdots + b_n = 0, \quad \xi''^n + b''_1\xi''^{n-1} + \cdots + b''_n = 0,
\end{equation}
so bezeichne man mit
\[\omega_1, \omega_2, \ldots, \omega_m\]
die Produkte
\[\xi'^i \xi''^k \quad (i = 0, \ldots, n'-1; \; k = 0, \ldots, n''-1)\]
und mit $\xi$ eine der drei Funktionen $\xi' + \xi''$, $\xi' - \xi''$, $\xi'\xi''$, dann kann man mit Hilfe der Gleichungen (6) setzen:
\begin{equation}
\xi\omega_i = \chi_{i,1}\omega_1 + \chi_{i,2}\omega_2 + \cdots + \chi_{i,m}\omega_m \quad (i = 1, 2, \ldots, m)
\end{equation}
worin $m = n'n''$ ist, und die $\chi_{i,k}$ ganze rationale Funktionen von $z$ sind. Aus (7) ergibt sich durch Elimination der $\omega_i$:
\[\begin{vmatrix}
\chi_{1,1} - \xi & \chi_{1,2} & \cdots & \chi_{1,m} \\
\chi_{2,1} & \chi_{2,2} - \xi & \cdots & \chi_{2,m} \\
\vdots & \vdots & \ddots & \vdots \\
\chi_{m,1} & \chi_{m,2} & \cdots & \chi_{m,m} - \xi
\end{vmatrix} = 0\]
und dies ist eine Gleichung für $\xi$ von der Form (5).

Durch wiederholte Anwendung dieses Satzes folgt, daß jede ganze rationale Funktion von ganzen Funktionen wieder eine ganze Funktion ist.

\begin{defin}
Eine ganze Funktion $\xi$ heißt durch eine ganze Funktion $\xi'$ \emph{teilbar}, wenn eine dritte ganze Funktion $\xi''$ existiert, so daß
\begin{equation}
\xi = \xi' \cdot \xi''
\end{equation}
ist.
\end{defin}

Aus dieser Definition ergibt sich ohne weiteres: Ist $\xi$ teilbar durch $\xi'$ und $\xi'$ teilbar durch $\xi''$, so ist auch $\xi$ durch $\xi''$ teilbar. Sind $\xi'$ und $\xi''$ durch $\xi$ teilbar, so ist auch $\xi' + \xi''$ durch $\xi$ teilbar.

Sind $\xi_1, \xi_2, \xi_3, \ldots$ durch $\xi$ teilbar und $\xi'_1, \xi'_2, \xi'_3, \ldots$ beliebige ganze Funktionen, so ist auch $\xi'_1\xi_1 + \xi'_2\xi_2 + \xi'_3\xi_3 + \cdots$ durch $\xi$ teilbar.

\newpage
\section*{\S 176. Minimalbasis und Körperdiskriminante.}
\addcontentsline{toc}{section}{\S 176. Minimalbasis und Körperdiskriminante.}

Da man jede Funktion des Körpers $\Omega$ durch Multiplikation mit einer ganzen rationalen Funktion in eine ganze algebraische Funktion von $z$ verwandeln kann, so gibt es auch Körperbasen, die aus ganzen Funktionen bestehen (z.\,B. nach der Bezeichnung in \S 171, (2) die Potenzen von $a_0\vartheta$).

Ist nun
\begin{equation}
\omega_1, \omega_2, \ldots, \omega_n
\end{equation}
eine solche aus ganzen Funktionen bestehende Basis, so ist jede in der Form
\begin{equation}
\xi = x_1\omega_1 + x_2\omega_2 + \cdots + x_n\omega_n
\end{equation}
enthaltene Funktion, wenn die $x_1, x_2, \ldots, x_n$ ganze rationale Funktionen sind, eine ganze Funktion in $z$. Es ist aber nicht gesagt, daß auch umgekehrt in der Form (2) mit ganzen Koeffizienten alle ganzen Funktionen in $\Omega$ enthalten sind.

Gibt es also ganze Funktionen in der Form (2), in der die $x_i$ nicht alle ganze rationale Funktionen sind, so können wir eine ganze Funktion finden:
\[\xi = x_1\omega_1 + x_2\omega_2 + \cdots + x_n\omega_n,\]
in der die $x_i$ ganz und nicht alle durch eine gewisse ganze rationale Funktion $z - c$ teilbar sind.

Reduzieren wir die $x_i$ auf ihre konstanten Reste (nach $z - c$), so ergibt sich eine ganze Funktion $\xi$ in der Form:
\begin{equation}
\xi - c_1\omega_1 - c_2\omega_2 - \cdots - c_n\omega_n = (z - c)\xi',
\end{equation}
in der die Konstanten $c_1, c_2, \ldots, c_n$ jedenfalls nicht alle verschwinden. Ist etwa $c_1$ von Null verschieden, so können wir $\omega_1$ durch $\xi, \omega_2, \ldots, \omega_n$ ausdrücken und erhalten eine neue ganzzahlige Basis von $\Omega$:
\begin{equation}
\xi, \omega_2, \ldots, \omega_n.
\end{equation}

Für die Diskriminante ergibt sich aber nach \S 173, (6) die Beziehung:
\begin{equation}
\Delta(\xi, \omega_2, \ldots, \omega_n) = (z - c)^2 \Delta(\omega_1, \omega_2, \ldots, \omega_n).
\end{equation}
Beide Diskriminanten sind ganze rationale Funktionen von $z$, aber die Diskriminante der Basis (4) ist von niedrigerem Grade als die Diskriminante der Basis (1).

Wenn wir auf diese Weise den Grad der Diskriminante zu erniedrigen fortfahren, müssen wir schließlich zu einer aus ganzen Funktionen bestehenden Basis $\beta_1, \beta_2, \ldots, \beta_n$ kommen der Eigenschaft, daß in der Form (2) mit ganzen rationalen $x_i$ alle und nur die ganzen Funktionen in $\Omega$ ausgedrückt sind.

\begin{defin}
Eine ganzzahlige Basis $\beta_1, \beta_2, \ldots, \beta_n$ des Körpers $\Omega$ heißt eine \emph{Minimalbasis}, wenn in der Form
\begin{equation}
\xi = x_1\beta_1 + x_2\beta_2 + \cdots + x_n\beta_n
\end{equation}
mit ganzen rationalen $x_1, x_2, \ldots, x_n$ alle ganzen Funktionen des Körpers $\Omega$ darstellbar sind.
\end{defin}

Eine solche Minimalbasis existiert also immer.

\begin{satz}
Ein aus einer Minimalbasis $\beta_1, \beta_2, \ldots, \beta_n$ abgeleitetes System ganzer Funktionen
\begin{equation}
\beta'_i = \chi_{i,1}\beta_1 + \chi_{i,2}\beta_2 + \cdots + \chi_{i,n}\beta_n
\end{equation}
ist dann und nur dann eine Minimalbasis, wenn die Determinante
\begin{equation}
|\chi_{i,k}| = c
\end{equation}
eine von Null verschiedene Konstante ist.
\end{satz}

Denn hat diese Determinante, die eine rationale Funktion von $z$ ist, einen Linearfaktor $z - c$, so kann man die Konstanten $c_1, c_2, \ldots, c_n$ so bestimmen, daß die $n$ ganzen rationalen Funktionen
\[c_1\chi_{1,\nu} + c_2\chi_{2,\nu} + \cdots + c_n\chi_{n,\nu}\]
durch $z - c$ teilbar werden, ohne daß alle $c_i$ verschwinden. Dann ist aber
\[\frac{c_1\beta'_1 + c_2\beta'_2 + \cdots + c_n\beta'_n}{z - c}\]
eine ganze Funktion, und $\beta'_1, \beta'_2, \ldots, \beta'_n$ keine Minimalbasis.

Für die Diskriminante erhält man aus (7) nach \S 175, (6):
\begin{equation}
\Delta(\beta'_1, \ldots, \beta'_n) = c^2 \Delta(\beta_1, \ldots, \beta_n).
\end{equation}

Daraus folgt:
\begin{satz}
Die Diskriminante einer Minimalbasis ist, von einem konstanten Faktor abgesehen, von der besonderen Wahl der Basis unabhängig.
\end{satz}

Die Diskriminante einer Minimalbasis hat unter allen Diskriminanten aus ganzen Funktionen gebildeter Basen den niedrigsten Grad (daher der Name Minimalbasis). Es ist eine durch den Körper selbst, abgesehen von einem konstanten Faktor, eindeutig bestimmte ganze rationale Funktion von $z$. Sie wird daher auch die \emph{Diskriminante des Körpers} $\Omega$ genannt und mit $D(\Omega)$ bezeichnet.

\newpage
\section*{Fünfundzwanzigster Abschnitt. Funktionale.}
\addcontentsline{toc}{section}{Fünfundzwanzigster Abschnitt. Funktionale.}

% §177
\section{Rationale Funktionale}
\addcontentsline{toc}{section}{\S 177. Rationale Funktionale.}

Wir übertragen nun den Begriff des Funktionals, der uns im 17. Abschnitt des zweiten Bandes zur Begründung der Theorie der algebraischen Zahlen gedient hat, auf die algebraischen Funktionen. Wir adjungieren also dem Körper $\Omega$ beliebige Variable und rechnen damit nach den Regeln der Buchstabenrechnung. Es entsteht so ein erweiterter Körper $\tilde{\Omega}$, dessen Elemente \emph{Funktionale} heißen. Der Körper $\Omega$ selbst heißt der Funktionalkörper. Die Variablen sind hier nicht im Sinne der Analysis als Zeichen für veränderliche Zahlen, sondern als bloße Rechnungssymbole aufzufassen, und sind wohl zu unterscheiden von den Variablen $z$, $\vartheta$ des Körpers $\Omega$. Wir wollen diese Hilfsvariablen die \emph{Funktionalvariablen} nennen.

Aus dem Körper $Z$ der rationalen Funktionen von $z$ entsteht durch Adjunktion der Variablen der Körper $\tilde{Z}$ der rationalen Funktionale.

Eine ganze rationale Funktion der Variablen:
\[\Phi(x, y, \ldots)\]
mit ganzen rationalen Funktionen von $z$ als Koeffizienten heißt eine ganze Funktion in $\tilde{Z}$. Der größte gemeinschaftliche Teiler der Koeffizienten von $\Phi$ heißt der \emph{Teiler} der Funktion, und die Funktion heißt \emph{primitiv}, wenn die Koeffizienten keinen gemeinschaftlichen Teiler haben. Die primitiven rationalen Funktionen und ihre Quotienten werden auch \emph{Einheiten} im Körper $\tilde{Z}$ genannt.

Die Quotienten ganzer Funktionen in $\tilde{Z}$ sind die Funktionale in $\tilde{Z}$. Jedes Funktional $A$ in $\tilde{Z}$ kann in die Form gesetzt werden:
\begin{equation}
A = \frac{\Phi_1}{\Phi_2} = a \cdot \frac{E_1}{E_2},
\end{equation}
worin $\Phi_1$, $\Phi_2$, $E_1$, $E_2$ ganze Funktionen in $\tilde{Z}$ sind, darunter die Einheiten $E_1$, $E_2$, und $E$ ist eine Einheit in $\tilde{Z}$, die sich als gebrochene Funktion darstellt. $a$ ist der Quotient der Teiler von $\Phi_1$ und $\Phi_2$, also eine ganze oder gebrochene rationale Funktion von $z$.

Die Funktionen $a$ und $E$ in der Formel (1) sind durch $A$ völlig bestimmt, abgesehen von konstanten Faktoren, die hier die Rolle der numerischen Einheiten spielen. Wir wollen $a$ die \glqq Absolute`` von $A$ nennen.

Es gilt dann der Satz:

\begin{satz}
Die Absolute eines Produktes ist gleich dem Produkt der Absoluten der Faktoren.
\end{satz}

und wir definieren:

\begin{defin}
Ein rationales Funktional heißt \emph{ganz}, wenn seine Absolute eine ganze Funktion von $z$ ist.
\end{defin}

Aus 1. folgt, daß das Produkt zweier ganzer Funktionale in $\tilde{Z}$ wieder ein ganzes Funktional ist. Dasselbe folgt auch für die Summe und Differenz zweier ganzer Funktionale $A_1$, $A_2$ in $\tilde{Z}$. Setzen wir nämlich:
\[A_1 = a_1 \frac{E_1}{E}, \quad A_2 = a_2 \frac{E_2}{E},\]
worin $a_1$, $a_2$ ganze rationale Funktionen von $z$ sind; $E$, $E_1$, $E_2$ ganze Einheiten, so ist die Absolute von $A_1 + A_2$ der Teiler der ganzen Funktion $a_1E_1 + a_2E_2$, also auch eine ganze Funktion von $z$.

Also gilt der Satz:

\begin{satz}
Summe, Differenz und Produkt zweier ganzer Funktionale in $\tilde{Z}$ sind wieder ganze rationale Funktionale.
\end{satz}

Alle Einheiten und deren Reziproken sind nach der Definition als ganz zu bezeichnen. Ist die Absolute eines ganzen rationalen Funktionals $A$ linear ($= z - c$), so heißt $A$ ein rationales Primfunktional oder ein Primfunktional in $\tilde{Z}$.

\begin{satz}
Jedes ganze rationale Funktional läßt sich in Primfaktoren zerlegen und zwar, von Einheitsfaktoren abgesehen, nur auf eine Weise.
\end{satz}

Ist ein Produkt von ganzen Funktionalen in $\tilde{Z}$ durch ein Primfunktional teilbar, so ist wenigstens einer der Faktoren durch dieses Primfunktional teilbar.

\newpage
\section*{\S 178. Funktionale des Körpers $\Omega$.}
\addcontentsline{toc}{section}{\S 178. Funktionale des Körpers $\Omega$.}

Ein Funktional des Körpers $\Omega$ ist ein Ausdruck von der Form:
\begin{equation}
\Xi = \xi_0 + \xi_1\vartheta + \xi_2\vartheta^2 + \cdots + \xi_{n-1}\vartheta^{n-1},
\end{equation}
wo $\vartheta$ die den Körper $\Omega$ erzeugende algebraische Funktion ist, und $\xi_0, \xi_1, \ldots, \xi_{n-1}$ rationale Funktionale sind. $\Xi$ ist nur dann $= 0$, wenn alle Koeffizienten $\xi_0, \xi_1, \ldots, \xi_{n-1}$ verschwinden.

Jedes Funktional, das in bezug auf $\vartheta$ von höherem als dem $(n-1)$ten Grad ist, kann durch den Rest der Division durch $f(\vartheta)$ ersetzt, also auf den $(n-1)$ten Grad reduziert werden, und der Quotient zweier Ausdrücke (1), in der der Nenner nicht $= 0$ ist, kann wie in \S 171, 1. auf die Form (1) gebracht werden.

Nimmt man eine beliebige Basis $\eta_1, \eta_2, \ldots, \eta_n$ des Körpers $\Omega$, so kann $\Xi$ auch in die Form
\begin{equation}
\Xi = \sum_{i=1}^n \zeta_i \eta_i
\end{equation}
gesetzt werden.

Wendet man dies auf die Produkte $\Xi\eta_i = \zeta_1\eta_1 + \zeta_2\eta_2 + \cdots + \zeta_n\eta_n$ an, worin die $\zeta_{i,k}$ rationale Funktionale sind, so ergibt sich für $\Xi$ eine Gleichung:
\begin{equation}
\begin{vmatrix}
\zeta_{1,1} - \Xi & \zeta_{1,2} & \cdots & \zeta_{1,n} \\
\zeta_{2,1} & \zeta_{2,2} - \Xi & \cdots & \zeta_{2,n} \\
\vdots & \vdots & \ddots & \vdots \\
\zeta_{n,1} & \zeta_{n,2} & \cdots & \zeta_{n,n} - \Xi
\end{vmatrix} = 0,
\end{equation}
also:

\begin{satz}
Jedes Funktional $\Xi$ in $\Omega$ genügt einer Gleichung
\begin{equation}
\varphi(\Xi) = \Xi^n + A_1\Xi^{n-1} + \cdots + A_{n-1}\Xi + A_n = 0,
\end{equation}
worin $\varphi = t^n + A_1t^{n-1} + \cdots + A_n$ eine rationale Funktion $n$ten Grades von $t$ ist, deren Koeffizienten $A_1, A_2, \ldots, A_n$ Funktionale des Körpers $\tilde{Z}$ sind.
\end{satz}

Wir definieren die Norm und die Spur des Funktionals $\Xi$ wie in \S 172:
\begin{equation}
N(\Xi) = \begin{vmatrix}
\zeta_{1,1} & \zeta_{1,2} & \cdots & \zeta_{1,n} \\
\zeta_{2,1} & \zeta_{2,2} & \cdots & \zeta_{2,n} \\
\vdots & \vdots & \ddots & \vdots \\
\zeta_{n,1} & \zeta_{n,2} & \cdots & \zeta_{n,n}
\end{vmatrix},
\end{equation}
\begin{equation}
S(\Xi) = \zeta_{1,1} + \zeta_{2,2} + \cdots + \zeta_{n,n},
\end{equation}
und die Sätze \S 172, (6) und (21) gelten unverändert auch von den Normen und Spuren der Funktionale.

Es ist dann, wie in \S 172, (8),
\begin{equation}
\varphi(t) = N(t - \Xi),
\end{equation}
und die Diskriminante eines Funktionalsystems ist wie in \S 173, (1) erklärt, und wenn das Funktional $\Xi$ in einem in $\Omega$ enthaltenen Körper $e$ten Grades $\Omega_1$ enthalten ist, so genügt $\Xi$ einer Gleichung $e$ten Grades $\varphi_1(t) = 0$. Wie in \S 172, (19) wird bewiesen, daß
\[N(t - \Xi)\]
eine Potenz von $\varphi_1(t)$ ist.

\newpage
\section*{\S 179. Ganze Funktionale des Körpers $\Omega$.}
\addcontentsline{toc}{section}{\S 179. Ganze Funktionale des Körpers $\Omega$.}

\begin{defin}
Ein Funktional $\Xi$ in $\Omega$ heißt \emph{ganz}, wenn es einer Gleichung
\[F(\Xi) = \Xi^n + A_1\Xi^{n-1} + \cdots + A_{n-1}\Xi + A_n = 0\]
genügt, deren Koeffizienten $A_1, \ldots, A_n$ ganze rationale Funktionale sind.
\end{defin}

Aus dem Gaussschen Satz (Bd. I, \S 2) folgt, da man nach \S 177 die Primfaktoren in $\tilde{Z}$ kennt, daß, wenn $F(t)$ in $\tilde{Z}$ reduzibel ist, auch jede in $F(t)$ aufgehende Funktion ganze Koeffizienten hat, insbesondere also auch die Gleichung niedrigsten Grades $\varphi(\Xi) = 0$, der $\Xi$ genügt.

Die notwendige und hinreichende Bedingung für ein ganzes Funktional $\Xi$ ist also die, daß
\[N(t - \Xi)\]
bei Ordnung nach Potenzen von $t$ ganze rationale Funktionale zu Koeffizienten hat.

Wie in \S 175 beweist man die Sätze:

\begin{satz}
Summe, Differenz, Produkt von ganzen Funktionalen in $\Omega$ sind wieder ganze Funktionale.
\end{satz}

\begin{satz}
Ist $\Xi$ ein ganzes Funktional in $\Omega$ nach der Definition 6. und zugleich rational, so ist es ein ganzes Funktional in $\tilde{Z}$ (nach der Definition 2.).
\end{satz}

\begin{satz}
Ist $\Xi$ ein ganzes Funktional in $\Omega$, so ist $N(\Xi)$ ein ganzes Funktional in $\tilde{Z}$, und die Absolute von $N(\Xi)$ heißt die \glqq absolute Norm von $\Xi$``.
\end{satz}

Die absolute Norm von $\Xi$ ist also eine Funktion in $Z$.

Bezeichnen wir die absolute Norm mit $N_a(\Xi)$, so gilt der Satz:
\begin{equation}
N_a(\Xi H) = N_a(\Xi) N_a(H).
\end{equation}

\newpage
\section*{\S 180. Teilbarkeit von Funktionalen. Einheiten.}
\addcontentsline{toc}{section}{\S 180. Teilbarkeit von Funktionalen. Einheiten.}

Es sind nun die Sätze von Bd. II, \S 155 mit ganz geringen Modifikationen, auf die im folgenden aufmerksam gemacht ist, zu wiederholen. Zur Vereinfachung des Ausdruckes sollen hier mit den kleinen griechischen Buchstaben ganze Funktionale in $\Omega$, mit den kleinen lateinischen Buchstaben ganze Funktionale in $\tilde{Z}$ bezeichnet sein. Dann haben wir:

\begin{defin}
Wenn $\beta$ von Null verschieden ist, so heißt $\alpha$ durch $\beta$ \emph{teilbar}, wenn $\alpha/\beta = \gamma$ ein ganzes Funktional in $\Omega$ ist.
\end{defin}

\begin{satz}
Sind $\xi$, $\eta$, $\ldots$ beliebige ganze Funktionale in $\Omega$ und $\alpha$, $\beta$, $\ldots$ durch $\delta$ teilbar, so ist auch $\xi\alpha + \eta\beta + \cdots$ durch $\delta$ teilbar.
\end{satz}

\begin{defin}
Ein ganzes Funktional $\varepsilon$, dessen Reziprokes $1/\varepsilon$ ganz ist, d.\,h. ein Teiler der Zahl $1$, heißt eine \emph{Einheit} in $\Omega$. Eine Einheit ist Teiler eines jeden ganzen Funktionals. Produkt und Quotient zweier Einheiten sind wieder Einheiten.
\end{defin}

\begin{defin}
Zwei ganze Funktionale $\alpha$, $\beta$, die gegenseitig durcheinander teilbar sind, heißen \emph{assoziiert}. Ihr Quotient $\alpha/\beta = \varepsilon$ ist eine Einheit. Zwei Funktionale $\alpha$ und $\alpha\varepsilon$ sind assoziiert.
\end{defin}

\begin{satz}
Ist $\alpha$ teilbar durch $\beta$, so ist jedes mit $\alpha$ assoziierte Funktional teilbar durch jedes mit $\beta$ assoziierte Funktional.
\end{satz}

\begin{satz}
Sind zwei ganze Funktionale mit einem dritten assoziiert, so sind sie auch untereinander assoziiert.
\end{satz}

\begin{satz}
Die Norm $N(\varepsilon)$ ist durch $\varepsilon$ teilbar.
\end{satz}
Dies folgt aus der Gleichung (6), \S 178: $\varphi(\varepsilon) = 0$, deren letzter Koeffizient $A_n = \pm N(\varepsilon)$ ist. Denn danach ist:
\[N(\varepsilon) = \varepsilon(\varepsilon^{n-1} + A_1\varepsilon^{n-2} + \cdots + A_{n-1}).\]

\begin{satz}
Ein ganzes Funktional, dessen absolute Norm eine Konstante ist, ist eine Einheit, und umgekehrt ist die absolute Norm einer Einheit $\varepsilon$ eine von Null verschiedene Konstante.
\end{satz}
Denn wenn $\varepsilon$ eine Einheit ist, so sind $\varepsilon$ und $1/\varepsilon$ ganze Funktionale in $\Omega$; folglich $N(\varepsilon)$ und $N(1/\varepsilon) = 1/N(\varepsilon)$ ganze Funktionale in $\tilde{Z}$, also $N(\varepsilon)$ eine Einheit in $\tilde{Z}$. Umgekehrt ist $N(\varepsilon)$ durch $\varepsilon$ teilbar, also, wenn $N(\varepsilon)$ eine Einheit in $\tilde{Z}$ ist, $\varepsilon$ ein Teiler von $1$, d.\,h. eine Einheit.

\begin{satz}
Es gibt ganze rationale Funktionen von $z$, z.\,B. die absolute Norm von $\alpha$, die durch $\alpha$ teilbar sind. Ist $\alpha$ eine durch $\alpha$ teilbare ganze rationale Funktion von $z$ von möglichst niedrigem Grade, so ist jede andere durch $\alpha$ teilbare ganze rationale Funktion von $z$ durch $a$ teilbar.
\end{satz}

\newpage
\section*{\S 181. Größter gemeinschaftlicher Teiler.}
\addcontentsline{toc}{section}{\S 181. Größter gemeinschaftlicher Teiler.}

Sind $\alpha$, $\beta$, $\ldots$ ganze Funktionale in $\Omega$ und $x$, $y$, $\ldots$ Variable, die in $\alpha$, $\beta$, $\ldots$ nicht vorkommen, so ist
\begin{equation}
\delta = \alpha x + \beta y + \cdots
\end{equation}
nach \S 180, 2. ebenfalls ein ganzes Funktional, und zwar ist $\delta$ teilbar durch jeden gemeinsamen Teiler von $\alpha$, $\beta$, $\ldots$ Sind $\xi_0$, $\eta_0$, $\ldots$ beliebige ganze Funktionen oder Funktionale in $\tilde{Z}$, so ist, wie jetzt bewiesen werden soll,
\begin{equation}
\frac{\alpha\xi_0 + \beta\eta_0 + \cdots}{\delta}
\end{equation}
durch $\delta$ teilbar.

Denn bezeichnen wir die absolute Norm von $\delta$ mit $D$, so ist
\begin{equation}
\frac{N(\delta x)}{D} = E(\xi, \eta, \ldots)
\end{equation}
und $E(\xi, \eta, \ldots)$ ist eine Einheit in $\tilde{Z}$, zugleich aber eine ganze homogene Funktion der Variablen $\xi$, $\eta$, $\ldots$.

Bedeutet $t$ eine neue Variable, so ist
\[N(\xi t - \delta_1) = DE(\xi t - X, \eta t - Y, \ldots);\]
also, wenn man nach absteigenden Potenzen von $t$ ordnet und den Faktor $N(\xi) = DE$ beiderseits forthebt:
\[\frac{\delta_1}{\delta} = G_0 + G_1 \frac{1}{t} + G_2 \frac{1}{t^2} + \cdots + G_m \frac{1}{t^m},\]
worin $G_0$, $G_1$, $\ldots$ keinen anderen Nenner als $E$ haben, und folglich ganze Funktionale in $\tilde{Z}$ sind.

Demnach ist nach der Definition \S 179, 6. auch $\delta_1:\delta$ ein ganzes Funktional, wie bewiesen werden sollte.

Demnach erhalten wir folgende Definitionen:

\begin{defin}
Das Funktional $\delta = \alpha x + \beta y + \cdots$ und jedes mit $\delta$ assoziierte Funktional ist der \emph{größte gemeinschaftliche Teiler} von $\alpha$, $\beta$, $\ldots$.
\end{defin}

\begin{defin}
Ist $\alpha x + \beta y$ eine Einheit, so heißen $\alpha$ und $\beta$ \emph{teilerfremd} oder \emph{relativ prim}. Gibt es Funktionale $\xi$, $\eta$, für die $\alpha\xi + \beta\eta$ eine Einheit ist, so sind $\alpha$, $\beta$ relativ prim.
\end{defin}

\begin{satz}
Ist $\alpha$ relativ prim zu $\beta$ und zu $\gamma$, so ist es auch relativ prim zu $\beta\gamma$.
\end{satz}
Denn nach Voraussetzung sind
\[\alpha\xi + \beta\eta = \varepsilon, \quad \alpha\mu + \gamma\nu = \varepsilon'\]
Einheiten ($\xi$, $\eta$, $\mu$, $\nu$ neue Variablen). Also
\[\alpha(\mu\varepsilon + \gamma\nu\alpha + \beta\eta\mu + \beta\eta\gamma\nu) + \beta\gamma\eta\nu = \varepsilon\varepsilon',\]
also nach 3. $\alpha$ und $\beta\gamma$ relativ prim.

\begin{satz}
Ist $\alpha$ relativ prim zu $\beta$ und $\alpha\mu$ durch $\beta$ teilbar, so ist $\mu$ durch $\beta$ teilbar.
\end{satz}
Denn aus der Voraussetzung folgt:
\[\alpha\xi + \beta\eta = \varepsilon, \quad \alpha\mu = \beta\nu,\]
woraus, da $\varepsilon$ eine Einheit ist, sich der Beweis ergibt.

\newpage
\section*{\S 182. Primfunktionale in $\Omega$.}
\addcontentsline{toc}{section}{\S 182. Primfunktionale in $\Omega$.}

\begin{defin}
Ein ganzes Funktional $\pi$ in $\Omega$ heißt ein \emph{Primfunktional}, wenn es keine Einheit ist und außer durch die Einheiten nur durch die mit ihm assoziierten Funktionale teilbar ist.
\end{defin}

Daraus folgt:

\begin{satz}
Ist ein Produkt $\alpha\beta$ durch $\pi$ teilbar, so ist wenigstens einer der beiden Faktoren, $\alpha$ oder $\beta$, durch $\pi$ teilbar.
\end{satz}
Denn ist weder $\alpha$ noch $\beta$ durch $\pi$ teilbar, so sind beide relativ prim zu $\pi$ und nach \S 181, 3. ist $\pi$ auch relativ prim zu $\alpha\beta$.

\begin{satz}
Jedes von Null verschiedene ganze Funktional $\alpha$ ist durch ein Primfunktional teilbar.
\end{satz}
Denn ist $\alpha$ nicht selbst ein Primfunktional, so ist es durch ein von $\alpha$ verschiedenes Funktional $\alpha'$ teilbar. Ist $\alpha = \alpha'\beta$, so ist $\beta$ keine Einheit und $N_a(\alpha')$ ist ein ganzer Teiler von $N_a(\alpha)$. Die hier vorkommenden absoluten Normen sind ganze Funktionen in $z$ und der Grad von $N_a(\alpha')$ ist kleiner als der Grad von $N_a(\alpha)$. Ist $\alpha'$ noch kein Primfunktional, so kann man so fortfahren und muß schließlich auf einen Primfaktor $\pi$ von $\alpha$ kommen.

In dieser Weise schließt man weiter wie in Bd. I, \S 158, wobei nur an Stelle der dort benutzten Größe der ganzen rationalen Zahlen hier die Höhe des Grades ganzer rationaler Funktionen der Variablen $z$ tritt. So erhält man auch die Sätze:

\begin{satz}
Jedes ganze Funktional $\alpha$ in $\Omega$, das keine Einheit ist, läßt sich in eine endliche Anzahl von Primfaktoren zerlegen, und zwar nur auf eine Weise, wenn assoziierte Funktionale als nicht verschieden betrachtet werden.
\end{satz}

Im folgenden müssen wir von dem Satz Gebrauch machen, daß eine ganze rationale Funktion von $z$ mit numerischen Koeffizienten in lineare Faktoren zerlegbar ist, mit anderen Worten, wir müssen den Fundamentalsatz der Algebra von der Wurzelexistenz voraussetzen. Es folgt daraus zunächst:

\begin{satz}
Die ganze rationale Funktion von $z$ niedrigsten Grades, die durch ein Primfunktional $\pi$ teilbar ist, ist eine lineare Funktion $z - c$.
\end{satz}

Nach \S 180, 9. gibt es überhaupt ganze rationale Funktionen von $z$, die durch $\pi$ teilbar sind. Zerlegt man eine solche Funktion in lineare Faktoren, so muß nach 2. einer dieser Linearfaktoren durch $\pi$ teilbar sein. Es können nicht zwei verschiedene Linearfunktionen $z - c$ und $z - c'$ durch dasselbe $\pi$ teilbar sein, weil sonst die von Null verschiedene Konstante $c' - c$ durch $\pi$ teilbar wäre.

\begin{satz}
Jede ganze Funktion $\alpha$ in $\Omega$ ist nach dem Modul $\pi$ mit einer Konstante $b$ kongruent, d.\,h. man kann die Konstante $b$ so wählen, daß $\alpha - b$ durch $\pi$ teilbar wird.
\end{satz}

Denn die Funktion $\alpha$ genügt nach \S 175 einer Gleichung:
\[\alpha^n + a_1\alpha^{n-1} + \cdots + a_{n-1}\alpha + a_n = 0,\]
worin die $a_1, \ldots, a_n$ ganze rationale Funktionen von $z$ sind. Sind $a'_1, \ldots, a'_n$ die Reste dieser Funktionen bei der Teilung durch $z - c$, also Konstanten, so folgt:
\[\alpha^n + a'_1\alpha^{n-1} + \cdots + a'_{n-1}\alpha + a'_n \equiv 0 \pmod{\pi},\]
und wenn man die Funktion auf der linken Seite in Linearfaktoren zerlegt:
\begin{equation}
(\alpha - b)(\alpha - b')(\alpha - b'')\cdots \equiv 0 \pmod{\pi}.
\end{equation}
Es muß also wenigstens einer dieser Linearfaktoren durch $\pi$ teilbar sein, was zu beweisen war.

Aus der nachgewiesenen einwertigen Zerlegbarkeit der Funktionale in Primfaktoren ergeben sich weitgehende Folgerungen, von denen die wichtigsten hier angeführt werden sollen.

Wenn in einem Funktional $g$ des Körpers $\Omega$ eine gewisse Anzahl der Hilfsvariablen $x$, $y$, $\ldots$, nur im Zähler vorkommen, so möge
\[g(x, y, \ldots)\]
eine \emph{holomorphe Funktion} von $\xi$, $\eta$, $\ldots$ heißen. Von denen gilt der Satz:

\begin{satz}
Eine holomorphe Funktion ist nur dann ein ganzes Funktional in $\Omega$, wenn die Koeffizienten der geordneten Funktion $g$ ganze Funktionale sind.
\end{satz}

Es genügt offenbar, den Satz für holomorphe Funktionen einer Variablen
\begin{equation}
g = \alpha_0\xi^m + \alpha_1\xi^{m-1} + \cdots + \alpha_{m-1}\xi + \alpha_m
\end{equation}
nachzuweisen, weil er daraus durch vollständige Induktion allgemein bewiesen werden kann.

Dieser Beweis ergibt sich aber aus dem Gaussschen Satz:

Haben zwei holomorphe Funktionen
\[\xi = \alpha_0\xi^m + \alpha_1\xi^{m-1} + \cdots + \alpha_m, \quad \eta = \beta_0\xi^n + \beta_1\xi^{n-1} + \cdots + \beta_n\]
ein Produkt
\[\xi\eta = \gamma_0\xi^{m+n} + \gamma_1\xi^{m+n-1} + \cdots + \gamma_{m+n},\]
dessen Koeffizienten ganze Funktionale sind, so können die $\gamma_0$, $\gamma_1$, $\ldots$, $\gamma_{m+n}$ nur dann alle durch ein Primfunktional $\pi$ teilbar sein, wenn entweder alle $\alpha_0, \alpha_1, \ldots, \alpha_m$, oder alle $\beta_0, \beta_1, \ldots, \beta_n$ durch $\pi$ teilbar sind, was ganz so bewiesen wird, wie in Bd. I, \S 2.

Wenn nun (2) ein ganzes Funktional ist, ohne daß $\alpha_0, \alpha_1, \ldots, \alpha_m$ ganz sind, so kann man ein ganzes Funktional $u$ und darin einen Primfaktor $\pi$ so bestimmen, daß
\begin{equation}
\iota g = u(\alpha_0\xi^m + \alpha_1\xi^{m-1} + \cdots + \alpha_m)
\end{equation}
durch $\pi$ teilbar ist, $\alpha_0, \alpha_1, \ldots, \alpha_m$ aber nicht alle durch $\pi$ teilbar sind.

Ist nun $\iota$ ganz, so genügt es einer Gleichung:
\begin{equation}
E\xi^m = A_1E\xi^{m-1} + A_2E\xi^{m-2} + \cdots + A_mE
\end{equation}
und daraus durch Multiplikation mit $u^m$:
\begin{equation}
E(u\xi)^m = u(A_1E(u\xi)^{m-1} + A_2E(u\xi)^{m-2} + \cdots + A_mE),
\end{equation}
worin die $A_1, A_2, \ldots$ ganze Funktionen in $\tilde{Z}$, die $E$, $E_1, \ldots, E_m$ Einheiten in $\tilde{Z}$ sind. Die Einheiten $E$, $E_1, \ldots, E_m$ können zwar $\pi$ noch enthalten, können aber holomorph angenommen werden.

Ordnen wir rechter Hand und linker Hand von (5) nach $\xi$, so sind die Koeffizienten von $E\xi^m$ nicht alle durch $\pi$ teilbar. Ebenso sind nach dem erwähnten Satz die Koeffizienten von $\iota\xi^m$ nicht alle durch $\pi$ teilbar, während auf der rechten Seite alle Koeffizienten den Faktor $u$, also auch den Faktor $\pi$ haben. Das ist unmöglich und damit unser Satz bewiesen (Bd. II, \S 159).

\begin{satz}
Ein holomorphes ganzes Funktional $g$ ist der größte gemeinschaftliche Teiler aller seiner Koeffizienten.
\end{satz}

Sind $\alpha$, $\beta$, $\ldots$ die Koeffizienten von $g$, so ist $g$ jedenfalls durch den größten gemeinschaftlichen Teiler von $\alpha$, $\beta$, $\ldots$ teilbar.

Zerlegen wir eine Linearfunktion $z - c$ in ihre funktionalen Primfaktoren, so können wir in diesen Primfunktionalen $\pi$ die Variablen beliebig bezeichnen, und da jedes Primfunktional in einem $z - c$ aufgehen muß, so können wir darin die Variablen von den Variablen $x$, $y$, $\ldots$ in $g$ verschieden annehmen.

Demnach können wir ein von den $x$, $y$, $\ldots$ freies, mit $g$ assoziiertes Funktional $\iota^{-1}g$ bestimmen. Da dann $g/\iota$ ganz ist, so müssen nach 7. die Koeffizienten von $g$ alle durch $\iota$, und folglich auch alle durch $g$ teilbar sein. Was zu beweisen war.

Damit ist zugleich bewiesen:

\begin{satz}
Ein holomorphes ganzes Funktional geht in ein assoziiertes über, wenn die Variablen anders bezeichnet werden.
\end{satz}

\begin{satz}
Jedes Funktional $\Xi$ in $\Omega$ geht durch Multiplikation mit einer Einheit in $\tilde{Z}$ in ein holomorphes Funktional über.
\end{satz}
Denn man kann zunächst $\Xi$ als Quotienten zweier ganzer holomorpher Funktionale
\[\Xi = \frac{\alpha}{\beta}\]
darstellen. Setzt man $N(\beta) = \beta\beta'$, so ist, wie die Formel (4) zeigt, in der $E_n = EN(\beta)$ ist, $\Xi\beta'$ mit einem holomorphen Funktional assoziiert, und es folgt:
\begin{equation}
E_n\Xi = \frac{\alpha'}{\beta}\beta' = \alpha',
\end{equation}
also
\begin{equation}
\Xi = \frac{\alpha'}{E_n}.
\end{equation}

\begin{satz}
Zu einem ganzen Funktional $\alpha$ kann man ein zweites $u$ so wählen, daß $u$ durch beliebig gegebene Primfunktionale $\pi_1$, $\pi_2$, $\ldots$ nicht teilbar wird, und daß
\begin{equation}
\alpha u = \Xi
\end{equation}
eine Funktion in $\Omega$ wird.
\end{satz}

Man stelle $\alpha$ nach 10. durch ein holomorphes Funktional $g$ dar. Dessen Koeffizienten sind nach 8. alle durch $g$, aber nicht alle durch $\pi_1$ teilbar. Es gibt also eine Funktion $\xi_1$, die durch $g$, aber nicht durch $\pi_1$ teilbar ist.

Nun setze man:
\[\xi = \xi_1 + \pi_1\xi_2 + \pi_1\pi_2\xi_3 + \cdots\]
und bestimme:
\begin{align*}
\xi_1 &\text{ teilbar durch } g, \text{ nicht teilbar durch } \pi_1, \\
\xi_2 &\text{ teilbar durch } g, \text{ nicht teilbar durch } \pi_2, \\
\xi_3 &\text{ teilbar durch } g, \text{ nicht teilbar durch } \pi_3, \\
&\vdots
\end{align*}
Dann genügt $\alpha u = \alpha\xi = \alpha\xi_1 + \alpha\pi_1\xi_2 + \alpha\pi_1\pi_2\xi_3 + \cdots$ den Forderungen des Satzes 11.

\begin{satz}
Jedes ganze Funktional $\alpha$ ist der größte gemeinschaftliche Teiler zweier Funktionen $\xi$, $\eta$ in $\Omega$.
\end{satz}

Um dies zu beweisen, nehmen wir $\xi = \alpha u$, teilbar durch $\alpha$, dann $\eta$ teilbar durch $\alpha$ und $\eta:u$ relativ prim zu $u$; dann ist $\alpha = \xi x + \eta y$ der größte gemeinschaftliche Teiler von $\xi$ und $\eta$.

\newpage
\section*{\S 183. Basen und Basisformen der Funktionale.}
\addcontentsline{toc}{section}{\S 183. Basen und Basisformen der Funktionale.}

Es sei jetzt
\begin{equation}
\beta_1, \beta_2, \ldots, \beta_n
\end{equation}
eine Minimalbasis des Körpers $\Omega$ (\S 176) und
\begin{equation}
\beta'_1, \beta'_2, \ldots, \beta'_n
\end{equation}
eine andere Basis von $\Omega$. Die lineare Substitution, die $\beta'$ mit dem $\beta$ zusammenhängen, sei:
\begin{equation}
\beta'_i = \alpha_{i,1}\beta_1 + \alpha_{i,2}\beta_2 + \cdots + \alpha_{i,n}\beta_n,
\end{equation}
worin die $\alpha_{i,k}$ ganze Funktionen in $\tilde{Z}$ sind, deren Determinante
\begin{equation}
A = |\alpha_{i,k}|
\end{equation}
nicht identisch verschwindet.

Sind $t_1$, $t_2$, $\ldots$, $t_n$ Variable, so ist
\begin{equation}
\Lambda = \alpha_{1,1}t_1 + \alpha_{2,1}t_2 + \cdots + \alpha_{n,1}t_n
\end{equation}
der größte gemeinschaftliche Teiler von $\xi_1$, $\xi_2$, $\ldots$, $\xi_n$. Setzt man für $t_1$, $t_2$, $\ldots$, $t_n$ ganze Funktionen in $\tilde{Z}$, so entstehen aus $\Lambda$ ganze Funktionen in $\Omega$, die alle durch $\Lambda$ teilbar sind.

Man nennt $\Lambda$ eine \emph{Basisform} und $\alpha_{1,1}$, $\alpha_{2,1}$, $\ldots$, $\alpha_{n,1}$ eine \emph{Basis} des Funktionals $\Lambda$, wenn man alle durch $\Lambda$ teilbaren Zahlen in $\Omega$ dadurch erhält, daß man für $t_i$ ganze Funktionen in $\tilde{Z}$ setzt.

Bilden die $\alpha_{i,1}$ eine solche Basis, so kann man ganze Funktionen $g_{i,k}$ in $\tilde{Z}$, so bestimmen, daß
\begin{equation}
g_{i,1}\alpha_{1,1} + g_{i,2}\alpha_{2,1} + \cdots + g_{i,n}\alpha_{n,1} = \beta_i
\end{equation}
wird.

Daraus folgt:
\begin{equation}
\sum_i g_{i,k}\alpha_{i,1} = \beta_k,
\end{equation}
wenn
\begin{equation}
u_k = \sum_i g_{i,k}t_i\end{equation}
gesetzt ist, und indem man für $\xi_k$ in (7) die Ausdrücke (3) substituiert:
\begin{equation}
\Lambda = \sum_k \nu_k \beta'_k.
\end{equation}

Nach der Definition der Norm (\S 172) ist also $N(\Lambda)$ die Determinante aus den Koeffizienten $(\sum_i \alpha_{i,1}g_{i,k})$, und diese kann man nach dem Multiplikationssatz der Determinanten
\begin{equation}
N(\Lambda) = A \cdot T
\end{equation}
zerlegen, worin $A$ die Bedeutung (4) hat, und $T$ die Determinante ist:
\begin{equation}
T = |g_{i,k}|.
\end{equation}

Hierin ist $T$ eine homogene Funktion $n$ten Grades der $t_i$, deren Koeffizienten ganze Funktionen in $\tilde{Z}$ sind. $A$ ist selbst eine ganze Funktion in $\Omega$. Man beweist nun wie in Bd. II, \S 164, daß $T$ eine Einheit ist.

Wäre das nämlich nicht der Fall, so hätte $T$ irgend einen Linearfaktor $z - c$, und man könnte nach einem elementaren Determinantensatz holomorphe ganze Funktionale $\gamma_i$ in $z$, die nicht alle durch $z - c$ teilbar sind, so bestimmen, daß
\[\gamma_1t_1 + \gamma_2t_2 + \cdots + \gamma_nt_n = 0\]
durch $z - c$ teilbar wird (man kann den $\gamma_i$ sogar konstante Koeffizienten geben, da man sie auf ihre Reste nach $z - c$ reduzieren kann).

Setzt man nun
\[\omega = \gamma_1\beta'_1 + \gamma_2\beta'_2 + \cdots + \gamma_n\beta'_n,\]
so folgt aus (7)
\[\Lambda\omega = \nu_1\xi_1 + \nu_2\xi_2 + \cdots + \nu_n\xi_n,\]
und daraus folgt, weil die $\xi_k$ durch $\Lambda$, die $\nu_k$ durch $z - c$ teilbar sind, daß $\omega$ durch $z - c$ teilbar ist, was der Definition der Minimalbasis widerspricht. Demnach ergibt sich aus (10), daß $A$ die absolute Norm des Funktionals $\Lambda$ ist:
\begin{equation}
N_a(\Lambda) = A \cdot E,
\end{equation}
worin $E$ eine Einheit ist.

Für ein Primfunktional $\pi$ können wir leicht eine Basis finden. Es können nicht alle Elemente $\beta_1, \beta_2, \ldots, \beta_n$ einer Minimalbasis durch $\pi$ teilbar sein, weil ja durch die Linearform
\begin{equation}
\beta_0 = x_1\beta_1 + x_2\beta_2 + \cdots + x_n\beta_n
\end{equation}
auch die Funktion \glqq 1`` darstellbar sein muß. Nehmen wir an, $\beta_1$ sei nicht durch $\pi$ teilbar. Nach \S 182, 6. gibt es Konstanten $c_1, c_2, \ldots, c_n$, deren erste unbeschadet der Allgemeinheit $= 1$ angenommen werden kann, die den Bedingungen:
\[\beta_1 \equiv 1, \quad \beta_2 \equiv c_2, \quad \ldots, \quad \beta_n \equiv c_n \pmod{\pi}\]
genügen. Wir setzen:
\begin{equation}
\begin{aligned}
\xi_1 &= (z - c)\beta_1, \\
\xi_2 &= \beta_2 - c_2\beta_1, \\
&\vdots \\
\xi_n &= \beta_n - c_n\beta_1.
\end{aligned}
\end{equation}

Daß dies eine Basis von $\Omega$ ist, ersieht man sofort, wenn man die Funktion (13) so darstellt:
\begin{equation}
\xi = x_1\xi_1 + x_2\xi_2 + \cdots + x_n\xi_n = (x_1(z-c) - x_2c_2 - \cdots - x_nc_n)\beta_1 + x_2\beta_2 + \cdots + x_n\beta_n;
\end{equation}
da $\xi_2$, $\ldots$, $\xi_n$ durch $\pi$ teilbar sind, so kann $\xi$ nur dann durch $\pi$ teilbar sein, wenn die ganze Funktion in $\tilde{Z}$:
\[x_1(z-c) - x_2c_2 - \cdots - x_nc_n\]
durch $(z - c)$ teilbar und folglich $\xi$ durch die Formel (15) darstellbar ist. Hiernach ergeben die Formeln (10) und (14):
\begin{equation}
N_a(\pi) = z - c,
\end{equation}
also den Satz:

\begin{satz}
Die absolute Norm eines Primfunktionals ist eine lineare Funktion von $z$.
\end{satz}

Dieser Satz bedeutet einen wesentlichen Unterschied zwischen den Theorien der Zahlen und der Funktionen. In der Zahlentheorie gibt es Primideale verschiedener Grade, deren Norm eine Potenz einer Primzahl ist. In der Theorie der Funktionen haben wir nur Primfunktionale ersten Grades. Dies rührt daher, daß wir vermöge des Fundamentalsatzes der Algebra jede ganze Funktion einer Variablen nach dem Modul $\pi$ in lineare Faktoren zerlegen können.

Die Zahl der Primfaktoren, in die $(z - c)$ zerfällt, ist hiernach $n$, da $N(z-c) = (z-c)^n$ ist. Fassen wir die untereinander gleichen Primfaktoren zusammen, so ist
\begin{equation}
z - c = \pi_1^{e_1}\pi_2^{e_2}\cdots\pi_h^{e_h} \cdot E,
\end{equation}
worin
\begin{equation}
e_1 + e_2 + \cdots + e_h = n
\end{equation}
und $E$ eine Einheit ist.

\newpage
\section*{\S 184. Basisform und Verzweigungsfunktional.}
\addcontentsline{toc}{section}{\S 184. Basisform und Verzweigungsfunktional.}

Unter der \emph{Basisform} des Körpers $\Omega$ wollen wir die Linearform verstehen:
\begin{equation}
r = \beta_1t_1 + \beta_2t_2 + \cdots + \beta_nt_n,
\end{equation}
in der $t_1$, $t_2$, $\ldots$, $t_n$ die Funktionalvariablen sind, und deren Koeffizienten $\beta_1, \beta_2, \ldots, \beta_n$ eine Minimalbasis von $\Omega$ bilden. Diese Form ist der größte gemeinschaftliche Teiler von $\beta_1, \beta_2, \ldots, \beta_n$ und folglich eine Einheit.

Aus $r$ kann man alle ganzen Funktionen in $\Omega$ ableiten, indem man für die Variablen $t_1$, $t_2$, $\ldots$, $t_n$ ganze Funktionen in $\tilde{Z}$ setzt.

Setzen wir
\begin{equation}
F(t) = N(t - r) = t^n + A_1t^{n-1} + \cdots + A_{n-1}t + A_n,
\end{equation}
so sind die $A_1, A_2, \ldots, A_n$ holomorphe ganze Funktionale in $\tilde{Z}$, und $r$ genügt der Gleichung $n$ten Grades:
\[F(r) = 0.\]

Ist $\pi$ ein Primfunktional und $N_a(\pi) = z - c$, so kann man nach \S 182, 6. die Konstanten $c_1, c_2, \ldots, c_n$ den Kongruenzen
\begin{equation}
\beta_1 \equiv c_1, \quad \beta_2 \equiv c_2, \quad \ldots, \quad \beta_n \equiv c_n \pmod{\pi}
\end{equation}
bestimmen, und wenn
\begin{equation}
r_c = c_1t_1 + c_2t_2 + \cdots + c_nt_n
\end{equation}
setzt, so ist $r - r_c$ ein konstantes Funktional, das durch $\pi$ teilbar ist. Es ist nun zu beweisen:

\begin{satz}
Das Primfunktional $\pi$ ist der größte gemeinschaftliche Teiler von $r - r_c$ und $z - c$.
\end{satz}

Zunächst ergibt sich aus der Definition \S 181, daß
\begin{equation}
r - r_c = t_1(\beta_1 - c_1) + t_2(\beta_2 - c_2) + \cdots + t_n(\beta_n - c_n)
\end{equation}
der größte gemeinschaftliche Teiler von $(\beta_1 - c_1)$, $(\beta_2 - c_2)$, $\ldots$, $(\beta_n - c_n)$ ist. Ist nun
\begin{equation}
\beta_i - c_i = \pi \cdot \gamma_i \quad (i = 1, 2, \ldots, n),
\end{equation}
so ist $r - r_c$ durch $\pi$ teilbar, und wenn daher $(z - c)$ durch eine erste Potenz von $\pi$ teilbar ist, so ist $r - r_c$ die erste Potenz von $\pi$ teilbar.

Ist $\pi'$ ein zweites Primfunktional, das auch mit $\pi$ identisch sein kann, und $z - c$ durch $\pi\pi'$ teilbar, und $r \equiv r'_c \pmod{\pi'}$, so ist $(r - r_c):(z - c)$ relativ prim zu $\pi$ (nach 3.) und wegen (11) ist $\Phi_1(r_c) = 0 \pmod{\pi'}$. Daraus schließt man
\[\Phi_1(t) = (t - r_c)\Phi_2(t)\]
zunächst für den Modul $\pi'$, dann aber auch, da die Funktionale in $\tilde{Z}$ liegen, auch für den Modul $(z - c)$. Wir haben also:
\begin{equation}
\Phi(t) = (t - r_c)(t - r'_c)\Phi_2(t) \pmod{z - c}.
\end{equation}

Daraus schließen wir auf folgenden wichtigen Satz:

\begin{satz}
Es sei nach \S 183, (17)
\begin{equation}
z - c = \pi_1^{e_1}\pi_2^{e_2}\cdots\pi_h^{e_h} \cdot E
\end{equation}
in $n$ (gleiche oder verschiedene) Primfaktoren zerlegt, und
\begin{equation}
r_1, r_2, \ldots, r_h
\end{equation}
seien die konstanten Funktionale, denen $r$ nach den Modulen
\begin{equation}
\pi_1, \pi_2, \ldots, \pi_h
\end{equation}
kongruent wird.

Es sei $\Phi(t)$ ein ganzes holomorphes Funktional in $\Omega$ vom Grade $m$ in bezug auf $t$, das der Bedingung
\begin{equation}
\Phi(r) = 0 \pmod{z - c}
\end{equation}
genügt. Dann ist:
\begin{equation}
\Phi(t) \equiv (t - r_1)(t - r_2)\cdots(t - r_h) \Psi(t) \pmod{z - c},
\end{equation}
worin $\Psi(t)$ ein ebensolches Funktional vom Grade $m - h$ ist.
\end{satz}

Es folgt daraus:

\begin{satz}
Ein Funktional $\Phi$, das der Bedingung (18) genügt, kann nicht von niedrigerem als $n$tem Grade sein.
\end{satz}

Und auf die Funktion $F(t)$ angewandt, folgt die Formel:
\begin{equation}
F(t) = N(t - r) \equiv (t - r_1)(t - r_2)\cdots(t - r_h) \pmod{z - c}.
\end{equation}

Die Kongruenz (20) können wir in folgender Weise durch eine Gleichung ersetzen: Wir fassen, wie in (17) des vorigen Paragraphen, in der Zerlegung von $z - c$ in Primfaktoren die gleichen Faktoren zusammen und setzen:
\begin{equation}
\pi_1^{e_1} = \zeta_1, \quad \pi_2^{e_2} = \zeta_2, \quad \ldots, \quad \pi_h^{e_h} = \zeta_h,
\end{equation}
dann ist nach (20):
\begin{equation}
F(t) = (t - r_1)(t - r_2)\cdots(t - r_h) + (z - c)R(t),
\end{equation}
worin $R(t)$ ein ganzes Funktional zunächst in $\Omega$ ist. Die Formel (22) selbst sagt aber, daß $(z-c)R(t)$ in $\tilde{Z}$ enthalten ist, und daß es eine ganze Funktion der Funktionalvariablen $t$, $t_1$, $t_2$, $\ldots$, $t_n$ ist.

Wir können also auch die Ableitung der Gleichung (12) in bezug auf $t$ bilden und erhalten:
\[F'(t) = \frac{\partial}{\partial t}[(t - r_1)(t - r_2)\cdots(t - r_h)] + (z - c)R'(t),\]
und hierin liegt der Beweis des folgenden Hauptsatzes, den man erhält, wenn man $t = r_i$ setzt:

\begin{satz}
$F'(r)$ ist durch keinen Primfaktor von $(z - c)$ teilbar, dessen Exponent $e = 1$ ist. Ist $e > 1$, so ist $F'(r)$ durch $\zeta^{e-1}$, aber nicht durch $\zeta^e$ teilbar. Das Gleiche gilt von den übrigen Primfaktoren von $(z - c)$ und folglich ist $F'(r)$ teilerfremd zu allen Primfaktoren von $(z - c)$ in $\tilde{Z}$, die nicht das Quadrat eines Primfaktors in $\Omega$ teilbar sind.
\end{satz}

Es gibt also nur eine endliche Anzahl von Primfaktoren $(z - c)$, die durch eine höhere als die erste Potenz eines Primfaktors in $\Omega$ teilbar sind, und $F'(r)$ ist, in Primfaktoren zerlegt, das Produkt aller Faktoren $\zeta^{e-1}$, wenn $\zeta^e$ die höchste in $N(\pi) = z - c$ aufgehende Potenz von $\zeta$ ist.

Das Funktional $F'(r)$ wird das \emph{Verzweigungsfunktional} des Körpers $\Omega$ genannt.

Für die absolute Norm des Verzweigungsfunktionals erhalten wir:
\begin{equation}
N_a(F'(r)) = \prod (z - c)^{e-1},
\end{equation}
worin sich das Produkt $\prod$ auf alle Linearfaktoren $(z - c)$ erstreckt, die durch eine höhere als die erste Potenz eines Primfaktors teilbar sind.

Wir können noch zeigen, daß diese absolute Norm nichts anderes ist, als die Körperdiskriminante.

Um dies zu beweisen, setzen wir:
\begin{equation}
r_i = u_{i,1}\beta_1 + u_{i,2}\beta_2 + \cdots + u_{i,n}\beta_n,
\end{equation}
worin die $u_{i,1}$, $u_{i,2}$, $\ldots$, $u_{i,n}$ homogene Funktionen ersten Grades von $t_1$, $t_2$, $\ldots$, $t_n$ sind, und die Determinante
\begin{equation}
U = \begin{vmatrix}
u_{1,1} & u_{1,2} & \cdots & u_{1,n} \\
u_{2,1} & u_{2,2} & \cdots & u_{2,n} \\
\vdots & \vdots & \ddots & \vdots \\
u_{n,1} & u_{n,2} & \cdots & u_{n,n}
\end{vmatrix}
\end{equation}
ist eine Einheit.

Denn wäre $U$ keine Einheit, so hätte sie wenigstens einen Linearfaktor $z - c$, und man könnte die ganzen rationalen holomorphen Funktionale $\gamma_0, \gamma_1, \ldots, \gamma_{n-1}$ so bestimmen, daß
\[u_{r,0}\gamma_0 + u_{r,1}\gamma_1 + \cdots + u_{r,n-1}\gamma_{n-1} = 0 \pmod{z-c}\]
wäre, ohne daß alle $\gamma_i$ durch $z - c$ teilbar sind. Dann würde sich aus (20) eine Kongruenz ergeben:
\[\gamma_0t_0 + \gamma_1t_1 + \cdots + \gamma_{n-1}t_{n-1} = 0 \pmod{z-c},\]
was nach dem Satz 5. nicht möglich ist.

Nach \S 174, (17) und \S 176, (9) ist
\begin{equation}
N_a(F'(r)) = \pm \Delta(\beta_1, \beta_2, \ldots, \beta_n) \cdot U^2
\end{equation}
also ist
\begin{equation}
N_a(F'(r)) = \text{const} \cdot D(\Omega)
\end{equation}
die Körperdiskriminante, die hiernach durch (24) in lineare Faktoren zerlegt ist.

Die gewonnenen Resultate benutzen wir noch zum Beweis des folgenden Satzes:

\begin{satz}
Ist eine ganze Funktion in $\Omega$
\begin{equation}
\omega = \xi_1\beta_1 + \xi_2\beta_2 + \cdots + \xi_n\beta_n
\end{equation}
durch jedes in $z - c$ aufgehende Primfunktional teilbar, so ist $S(\omega)$ durch $z - c$ teilbar.
\end{satz}

Die $\xi_1, \xi_2, \ldots, \xi_n$ sind hier ganze Funktionen in $\tilde{Z}$. Die Funktion $\omega$ geht aus $r$ hervor durch die Substitution
\begin{equation}
t_1 = \xi_1, \quad t_2 = \xi_2, \quad \ldots, \quad t_n = \xi_n.
\end{equation}
Ist also $\omega$ durch $\pi$ teilbar, und $c_1, \ldots, c_n$ die den Kongruenzen (4) entsprechenden Werte, so ist
\[\xi_1c_1 + \xi_2c_2 + \cdots + \xi_nc_n = 0 \pmod{\pi}\]
und folglich auch, als Funktion in $\tilde{Z}$, durch $z - c$ teilbar.

Sind also $r_1, r_2, \ldots, r_m$ die konstanten Funktionale, denen $r$ nach den Primfaktoren $\pi_1, \pi_2, \ldots, \pi_m$ von $(z - c)$ kongruent ist, und gehen diese durch die Substitution (30) in $z_1, z_2, \ldots, z_m$ über, so sind alle diese Funktionen durch $(z - c)$ teilbar [weil sie n.\,V. durch einen Primteiler von $(z - c)$ teilbar sind]. Nun ist nach (22)
\begin{equation}
S(\omega) = \xi_1T_1 + \xi_2T_2 + \cdots + \xi_mT_m;
\end{equation}
und mithin
\[-S(\omega) \equiv z_1 + z_2 + \cdots + z_m = 0 \pmod{z-c},\]
wie zu beweisen war.

\newpage
\section*{\S 185. Die gebrochenen Funktionen in $z$ und die Taylorsche Entwickelung.}
\addcontentsline{toc}{section}{\S 185. Die gebrochenen Funktionen in $z$ und die Taylorsche Entwickelung.}

Eine gebrochene Funktion $\eta$ im Körper $\Omega$ kann auf unendlich viele Arten durch Multiplikation mit einer ganzen Funktion $v$ in eine ganze Funktion $u$ verwandelt werden. Es ist dann
\begin{equation}
\eta = \frac{u}{v},
\end{equation}
und $u$ heißt der Zähler, $v$ der Nenner von $\eta$. Diese beiden Funktionen sind durch $\eta$ nicht vollständig bestimmt, wohl aber sind die Funktionale bestimmt, die übrig bleiben, wenn alle gemeinschaftlichen Funktionalfaktoren im Zähler und Nenner herausgehoben werden.

Denn ist
\begin{equation}
\eta = \frac{u}{v} = \frac{u'}{v'},
\end{equation}
so muß jedes Primfunktional, das in $v$, aber nicht in $u$ aufgeht, in $v'$ aufgehen. Die so aus $\eta$ bestimmten Funktionale $\alpha$ und $\beta$ nennen wir Zählerfunktional und Nennerfunktional von $\eta$. Um $\eta$ durch Funktionen darzustellen, nehme man zunächst eine Funktion $v$, teilbar durch $\beta$, aber sonst beliebig und setze $v = \tau$; darin kann $\tau$ relativ prim zu einem beliebigen Funktional angenommen werden. Dann ist $\eta v$ eine durch $\tau$ teilbare ganze Funktion, und man setze
\[\eta v = \xi.\]
Das Funktional $\xi$ kann aber, wie wir später noch nachweisen werden, bei gegebenem $\beta$ nicht mehr beliebig sein.

Die Funktionale $\alpha$, $\beta$ werden wir auch kurz den Zähler und den Nenner von $\eta$ nennen.

Ist $\pi$ ein Primfunktional, so gibt es Zahlen (Konstanten) $a$, $b$, $a'$, $b'$, die den Kongruenzen
\[u \equiv a, \quad v \equiv b, \quad u' \equiv a', \quad v' \equiv b' \pmod{\pi}\]
genügen, und aus (2) folgt
\[ab' = a'b,\]
vorausgesetzt, daß $v$ und $v'$ nicht durch $\pi$ teilbar sind. Es ist dann $\eta - \frac{a}{b}$ eine Funktion in $\Omega$, in der der Zähler, aber nicht der Nenner durch $\pi$ teilbar ist. Ist $\omega$ eine ganze oder gebrochene Funktion, in der der Zähler durch $\pi$, aber nicht durch $\pi^2$, der Nenner auch nicht durch $\pi$ teilbar ist, so ist $(\eta - c):\omega$ eine Funktion, deren Nenner nicht durch $\pi$ teilbar ist und deren Zähler den Faktor $\pi$ einmal weniger enthält als $(\eta - c)$. Wir setzen:
\[\eta = c + \omega\eta_1.\]
Dieselbe Betrachtung wenden wir auf $\eta_1$ an und bekommen so eine beliebig fortzusetzende Reihe von Gleichungen:
\[\eta_1 = c_1 + \omega\eta_2, \quad \eta_2 = c_2 + \omega\eta_3, \quad \ldots\]
woraus sich ergibt:
\begin{equation}
\eta = c + c_1\omega + c_2\omega^2 + \cdots + c_{k-1}\omega^{k-1} + \omega^k\eta_k.
\end{equation}

Die Reihe der Konstanten $c$, $c_1$, $c_2$, $\ldots$, $c_{k-1}$ ist durch $\eta$ vollständig bestimmt und $\eta_1$, $\eta_2$, $\ldots$, $\eta_k$ ist eine Reihe von Funktionen in $\Omega$, deren Nenner nicht durch $\pi$ teilbar ist.

Man könnte diesen Ausdruck die \emph{Taylorsche Entwickelung} der Funktion $\eta$ nach Potenzen von $\omega$ nennen. Eine ähnliche Entwickelung läßt sich auch für eine Funktion $\eta$ geben, deren Nenner durch $\pi$ teilbar ist, nur daß dabei auch negative Potenzen auftreten: Man kann nämlich eine Potenz $\omega^m$ so bestimmen, daß der Nenner von $\omega^m\eta$ nicht mehr durch $\pi$ teilbar ist, und wenn man auf dieses Produkt die Entwickelung (4) anwendet, so folgt:
\begin{equation}
\eta = c_{-m}\omega^{-m} + c_{-m+1}\omega^{-m+1} + \cdots + c_{-1}\omega^{-1} + c_0 + c_1\omega + \cdots
\end{equation}

\newpage
\section*{\S 186. Birationale Transformation.}
\addcontentsline{toc}{section}{\S 186. Birationale Transformation.}

Wir kommen jetzt zu einem Gegenstand, der das Gebiet der algebraischen Funktionen vorzugsweise von dem der algebraischen Zahlen unterscheidet, und der Grund dafür ist, daß die Theorie der ganzen Funktionen und Funktionale für die Funktionen bei weitem nicht den Charakter der Invarianz hat, wie bei den Zahlen; das ist die \emph{birationale Transformation}.

Ist $z_1$ irgend eine nicht konstante Funktion des Körpers $\Omega$ und $N(t - z_1)$ die $f$te Potenz einer irreduziblen Funktion $e$ten Grades (\S 172, 6.), so besteht zwischen $z$ und $z_1$ eine Gleichung, die in bezug auf $z_1$ vom $e$ten Grade ist, deren Grad in bezug auf $z$, wenn Nenner und überflüssige Faktoren weggeschafft sind, mit $e_1$ bezeichnet werde. Diese Gleichung sei
\begin{equation}
g(z, z_1) = 0.
\end{equation}

Unter den verschiedenen Gleichungen dieser Form gibt es nur eine irreduzible und diese ist sowohl in bezug auf $z$ als in bezug auf $z_1$ von möglichst niedrigem Grade (Bd. I, \S 20).

Ist $\vartheta$ eine primitive Funktion des Körpers $\Omega$, so ist nach \S 172, (12)
\begin{equation}
1, \vartheta, \vartheta^2, \ldots, \vartheta^{f-1}
\end{equation}
eine Basis des Körpers $\Omega$, und folglich kann jede Funktion $\eta$ in $\Omega$ in der Form dargestellt werden:
\begin{equation}
\eta = \xi_0 + \xi_1\vartheta + \xi_2\vartheta^2 + \cdots + \xi_{f-1}\vartheta^{f-1},
\end{equation}
worin die $\xi_0, \xi_1, \ldots, \xi_{f-1}$ rationale Funktionen von $z$ und $z_1$ sind, also Zahlen des durch (1) bestimmten Körpers.

Setzen wir
\begin{equation}
\eta = \xi\vartheta^k,
\end{equation}
so lassen sich die Funktionen $\xi$ in jeder der beiden Formen darstellen:
\begin{equation}
\xi_i = \sum_{j=0}^{e-1} x_{i,j}z_1^j = \sum_{k=0}^{e_1-1} y_{i,k}z^k,
\end{equation}
worin die $x$ rationale Funktionen von $z$, die $y$ rationale Funktionen von $z_1$ sind.

\begin{satz}
Demnach läßt sich jede Funktion $\eta$ linear mit rationalen Koeffizienten in $z_1$ darstellen durch die $n_1$-Funktionen:
\begin{equation}
\eta_i^{(k)} = z^i\vartheta^k \quad (i = 0, 1, \ldots, e_1-1; \; k = 0, 1, \ldots, f-1),
\end{equation}
die wir in irgend einer Reihenfolge mit
\begin{equation}
\eta_1, \eta_2, \ldots, \eta_{n_1}
\end{equation}
bezeichnen, und diese Funktionen sind linear unabhängig, d.\,h. es besteht zwischen ihnen keine lineare Relation mit rationalen Koeffizienten in $z_1$:
\[\xi_1\eta_1 + \xi_2\eta_2 + \cdots + \xi_{n_1}\eta_{n_1} = 0,\]
außer wenn die $\xi_1, \xi_2, \ldots, \xi_{n_1}$ alle verschwinden.
\end{satz}

Denn $f$ ist der niedrigste Grad einer Gleichung für $\vartheta$ mit rationalen Koeffizienten in $z$ und $z_1$ (\S 172).

Betrachten wir nun das Funktional
\begin{equation}
\Upsilon = \xi_1\eta_1 + \xi_2\eta_2 + \cdots + \xi_{n_1}\eta_{n_1}
\end{equation}
mit den Funktionalvariablen $\xi_1, \xi_2, \ldots, \xi_{n_1}$, so ergibt sich ein System von Gleichungen:
\begin{equation}
\Upsilon^r = \chi_{1,r}\eta_1 + \chi_{2,r}\eta_2 + \cdots + \chi_{n_1,r}\eta_{n_1},
\end{equation}
worin $r$ jede natürliche Zahl, auch $0$, sein kann, und die $\chi_{i,r}$ rationale holomorphe Funktionale in $z_1$ sind.

Es ergibt sich daraus durch Elimination der $\eta$ eine Gleichung
\begin{equation}
D(\Upsilon) = \Upsilon^{n_1} + A_1\Upsilon^{n_1-1} + \cdots + A_{n_1-1}\Upsilon + A_{n_1} = 0
\end{equation}
höchstens vom Grade $n_1$, mit rationalen Funktionalen in $z_1$ als Koeffizienten.

Es ist zu beweisen, daß dieser Grad nicht kleiner als $n_1$ sein kann, oder was damit gleichbedeutend ist, daß die $\eta_1$, $\eta_2$, $\ldots$, $\eta_{n_1}$ rational durch $\Upsilon$ darstellbar sind. Dies ergibt sich, wenn man (10) nach einem der $\xi_i$ partiell differentiiert:
\[D'(\Upsilon)\eta_i + \frac{\partial A_1}{\partial\xi_i}\Upsilon^{n_1-1} + \cdots + \frac{\partial A_{n_1}}{\partial\xi_i} = 0.\]
Da nun $D'(\Upsilon)$ nicht verschwindet, wenn $m$ der möglichst niedrige Grad der Gleichung (10) ist, so erhält man daraus unmittelbar die gesuchte Darstellung von $\eta_i$. Also muß $m = n_1$ sein.

Damit ist bewiesen:

\begin{satz}
Das Funktional $\Upsilon$ genügt einer irreduziblen Gleichung $n_1$ten Grades in $z_1$.
\end{satz}

Setzen wir also in (9) $r = 0, 1, 2, \ldots, n_1-1$, so ergibt sich ein System linearer Gleichungen für $\eta_1, \eta_2, \ldots, \eta_{n_1}$, deren Determinante ein nicht verschwindendes Funktional in $z_1$ ist, und man kann also für die Variable $\xi_i$ solche konstante Werte setzen, daß diese Determinante auch dann nicht verschwindet.

Geht $\Upsilon$ durch diese Substitution in $\Upsilon_0$ über, so kann man aus dem System (9) die $\eta_1, \eta_2, \ldots, \eta_{n_1}$ rational durch $z_1$ und $\Upsilon_0$ ausdrücken.

Damit ist aber bewiesen:

\begin{satz}
Ist $z_1$ eine beliebige, nicht konstante Funktion in $\Omega$, und $\vartheta_1$ eine zweite Funktion desselben Körpers, die keiner Gleichung von niedrigerem als $n_1$ten Grade in $z_1$ genügt, so sind alle Funktionen des Körpers $\Omega$, also auch $z$ und $\vartheta$, rational durch $z_1$ und $\vartheta_1$ ausdrückbar.
\end{satz}

Zwischen irgend zwei Funktionen $\alpha$, $\beta$ des Körpers $\Omega$ besteht immer eine Gleichung mit konstanten Koeffizienten
\[G(\alpha, \beta) = 0,\]
und unter allen möglichen Gleichungen dieser Form ist eine von möglichst niedrigem Grade, sowohl in bezug auf $\alpha$ als in bezug auf $\beta$.

\newpage
\section*{Sechsundzwanzigster Abschnitt. Zahlenwerte der algebraischen Funktionen.}
\addcontentsline{toc}{section}{Sechsundzwanzigster Abschnitt. Zahlenwerte der algebraischen Funktionen.}

% §187
\section{Der Punkt}
\addcontentsline{toc}{section}{\S 187. Der Punkt.}

Es entsteht nun die Frage, wie man den algebraischen Funktionen eine Bedeutung im Gebiete der Zahlen beilegen kann. Wir setzen dabei das Gebiet der reellen und imaginären rationalen und irrationalen Zahlen als gegeben voraus, und in diesem Gebiete das Rechnen mit den vier Spezies. Wir erweitern dieses Gebiet durch Hinzufügung eines Zahlzeichens \glqq Unendlich``: $\infty$, und rechnen auch mit diesem Zeichen in bekannter Weise, so daß
\begin{equation}
\infty + a = \infty, \quad a \cdot \infty = \infty \quad (a \neq 0), \quad \infty \cdot \infty = \infty, \quad a : 0 = \infty \quad (a \neq 0), \quad a : \infty = 0
\end{equation}
ist, während den Symbolen $\infty \pm \infty$, $0 \cdot \infty$, $0:0$, $\infty:\infty$ keine Bedeutung zukommt (oder auch, in jedem einzelnen Falle nach Bedarf, ein beliebiger Zahlwert beigelegt wird).

Der Körper $\Omega$ besteht jetzt aus allen möglichen algebraischen Funktionen $\alpha$, $\beta$, $\gamma$, $\ldots$, die nach \S 186 als rationale Funktionen von zweien unter ihnen mit Zahlenkoeffizienten dargestellt werden können. Auf die Art dieser Darstellung kommt es zunächst nicht an.

Wir stellen nun folgende Definition auf:

\begin{defin}
Wenn alle Individuen $\alpha$, $\beta$, $\gamma$, $\ldots$ des Körpers $\Omega$ durch bestimmte Zahlwerte $\alpha_0$, $\beta_0$, $\gamma_0$, $\ldots$ in der Weise ersetzt werden, daß:
\begin{enumerate}[label=\arabic*.]
\item $\alpha = \alpha_0$, wenn $\alpha$ konstant,
\item $(\alpha + \beta)_0 = \alpha_0 + \beta_0$,
\item $(\alpha \cdot \beta)_0 = \alpha_0 \cdot \beta_0$,
\item $(\alpha : \beta)_0 = \alpha_0 : \beta_0$,
\end{enumerate}
so ordnen wir einem solchen Zusammentreffen von Werten einen Punkt $\mathfrak{P}$ zu. Wir sagen: $\alpha = \alpha_0$, oder $\alpha$ hat den Wert $\alpha_0$ im Punkte $\mathfrak{P}$. Zwei Punkte $\mathfrak{P}$ und $\mathfrak{P}'$ heißen dann und nur dann verschieden, wenn es wenigstens eine Funktion in $\Omega$ gibt, die in $\mathfrak{P}$ und $\mathfrak{P}'$ verschiedene Werte hat.
\end{defin}

Die Regeln 2. bis 5. versagen in den Ausnahmefällen, nämlich 2. und 3., wenn $\alpha_0$ und $\beta_0$ beide $\infty$ sind; 4., wenn $\alpha_0 = \beta_0 = \infty$ oder $\beta_0 = 0$ wird, und 5., wenn $\alpha_0$ und $\beta_0$ beide $= 0$ oder beide $= \infty$ sind.

Trotzdem müssen die Funktionen $(\alpha + \beta)_0$, $(\alpha\beta)_0$ und $(\alpha:\beta)_0$ auch in diesen Fällen bestimmte Werte haben, die aber nicht unmittelbar aus den Vorschriften 2. bis 5., sondern auf indirektem Wege bestimmt werden.

Der Punkt ist hiernach ein zu dem Körper $\Omega$ gehöriger invarianter Begriff, der nichts mit dem zufälligen Umstand zu tun hat, welche der Funktionen von $\Omega$ wir als die unabhängige Variable betrachten.

Um alle Punkte zu finden, verfahre man so: Ist $\mathfrak{P}$ ein Punkt, so gibt es Funktionen, die in ihm einen endlichen Wert haben, denn wenn eine Funktion $\xi$ in $\mathfrak{P}$ unendlich ist, so ist $1:\xi$ endlich. Eine solche Funktion nehme man als unabhängige Variable $z$ und bezeichne ihren Wert im Punkt $\mathfrak{P}$ mit $c$. Alle ganzen Funktionen von $z$ sind dann gleichfalls in $\mathfrak{P}$ endlich, denn ist $\xi$ eine ganze Funktion, die der Gleichung
\[\xi^n + a_1\xi^{n-1} + \cdots + a_{n-1}\xi + a_n = 0\]
genügt, so muß
\[\xi_0^n + a_1(c)\xi_0^{n-1} + \cdots + a_{n-1}(c)\xi_0 + a_n(c) = 0\]
in $\mathfrak{P}$ befriedigt sein, und da die $a_1, a_2, \ldots, a_n$ als ganze rationale Funktionen von $z$ in $\mathfrak{P}$ endlich sind, so kann $\xi_0$ nicht $= \infty$ sein.

Jede Kongruenz nach dem Modul $z - c$ muß daher im Punkt $\mathfrak{P}$ in eine richtige Gleichung übergehen, auch wenn diese Kongruenz zwischen holomorphen ganzen Funktionalen besteht. Ist also $r$ die Basisform von $\Omega$ in $\tilde{Z}$, so muß nach \S 184, (20) die Gleichung bestehen:
\[(t - r_1)(t - r_2)\cdots(t - r_h) = 0.\]
Es muß also einer der Faktoren der linken Seite in $\mathfrak{P}$ verschwinden. Ist dies $t - r_i$, so entspricht diesem Faktor ein Primfaktor $\pi$ von $z - c$. Es wird also in $\mathfrak{P}$
\begin{equation}
\beta_1 = c_1, \quad \beta_2 = c_2, \quad \ldots, \quad \beta_n = c_n,
\end{equation}
wenn $c_1, c_2, \ldots, c_n$ die durch die Kongruenzen
\begin{equation}
\beta_1 \equiv c_1, \quad \beta_2 \equiv c_2, \quad \ldots, \quad \beta_n \equiv c_n \pmod{\pi}
\end{equation}
bestimmten Zahlwerte sind. Hierdurch sind die Werte aller ganzen Funktionen von $z$ im Punkt $\mathfrak{P}$ bestimmt. Jede ganze Funktion $\omega$ ist nach dem Modul $z - c$ einem Ausdruck
\[\alpha_1\beta_1 + \alpha_2\beta_2 + \cdots + \alpha_n\beta_n\]
mit konstantem Koeffizienten $\alpha_i$ kongruent und ist dann und nur dann durch $\pi$ teilbar, wenn
\[\alpha_1c_1 + \alpha_2c_2 + \cdots + \alpha_nc_n = 0\]
ist. Demnach folgt:

\begin{satz}
Unter den ganzen Funktionen von $z$ haben nur die durch $\pi$ teilbaren in $\mathfrak{P}$ den Wert $0$.
\end{satz}

Da man jede gebrochene Funktion in $\Omega$ so darstellen kann, daß Zähler und Nenner nicht zugleich durch $\pi$ teilbar sind, so ist hierdurch der Wert einer jeden Funktion in $\mathfrak{P}$ bestimmt.

\begin{defin}
Der Primfaktor $\pi$ heißt \emph{durch den Punkt $\mathfrak{P}$ erzeugt}. Er kann nicht durch einen von $\mathfrak{P}$ verschiedenen Punkt erzeugt werden.
\end{defin}

Man kann auch umgekehrt aus jedem Primfaktor $\pi$ in $z$ einen Punkt erzeugen. Ist $z - c$ die durch $\pi$ teilbare Linearfunktion, so gebe man der Funktion $z$ den Wert $c$. Dann nehme man eine durch $\pi$, aber nicht durch $\pi^2$ teilbare Funktion $\omega$ und setze nach \S 185, (5) für jede Funktion $\eta$:
\begin{equation}
\eta = a\omega^m + \omega^{m+1}\eta_1,
\end{equation}
worin $m$ eine ganze rationale Zahl, $\eta_1$ eine Funktion, deren Nenner nicht durch $\pi$ teilbar ist. Der Funktion $\omega$ erteile man den Wert $0$ und der Funktion $\eta$ den Wert
\[
a, \text{ wenn } m = 0, \quad
0, \text{ wenn } m > 0, \quad
\infty, \text{ wenn } m < 0.
\]

\begin{satz}
Die so bestimmten Werte aller Funktionen $\eta$ genügen den Bedingungen 1. bis 5. und konstituieren also einen Punkt, der durch das Primfunktional $\pi$ erzeugt heißt.
\end{satz}

Man kann also aus einer einzigen Funktion $z$ und den Primfaktoren der Linearformen $z - c$ alle Punkte ableiten, in denen $z$ einen endlichen Wert hat. Um auch die übrigen zu bestimmen, muß man noch eine zweite Funktion, etwa $z_1 = 1:z$, zu Hilfe nehmen, die in den noch fehlenden, in endlicher Anzahl vorhandenen Punkten den endlichen Wert $0$ hat.

Die Gesamtheit der Punkte bilden die \emph{absolute Riemannsche Fläche}\footnote{Nach Riemanns Theorie der algebraischen Funktionen entspricht jeder Punkt der geschlossenen mehrblätterigen Fläche einem Punkte in unserem Sinne. Ist $z$ unabhängige Variable, so ist die Riemannsche Fläche $n$-blättrig über die $z$-Ebene ausgebreitet. Als absolute Riemannsche Fläche kann man irgend eine Fläche betrachten, auf die die Gesamtheit jener mehrblätterigen Flächen eindeutig und stetig bezogen werden kann.}.

\newpage
\section*{\S 188. Ordnungszahlen.}
\addcontentsline{toc}{section}{\S 188. Ordnungszahlen.}

\begin{defin}
Ist $\mathfrak{P}$ ein Punkt, so schreiben wir jeder Funktion $\eta$, die in diesem Punkte einen endlichen und von Null verschiedenen Wert $\eta_0$ erhält, in $\mathfrak{P}$ die Ordnung $0$ zu.
\end{defin}

\begin{defin}
Wenn $\eta$ die Gesamtheit der in $\mathfrak{P}$ verschwindenden Funktionen durchläuft, und $\omega$ eine von diesen Funktionen ist, für die alle $\eta:\omega$ endliche Werte haben, so hat $\omega$ die Ordnung $1$.
\end{defin}

Man sagt auch, $\omega$ wird in $\mathfrak{P}$ unendlich klein in der ersten Ordnung. Ist $\omega$ weder Null noch Unendlich, so hat $\omega'$ hiernach gleichfalls die Ordnung $1$. Umgekehrt ist, wenn $\omega$ und $\omega'$ beide die Ordnungszahl $1$ haben, $\omega:\omega'$ und $\omega':\omega$ weder Null noch Unendlich.

\begin{defin}
Ist $\omega$ von der ersten Ordnung und $\xi$ irgend eine Funktion in $\Omega$, so hat $\xi$ die Ordnungszahl $n$, wenn $\xi\omega^{-n}$ in $\mathfrak{P}$ weder Null noch unendlich wird.
\end{defin}

Diese Bestimmung der Ordnung $n$ ist unabhängig davon, welche Funktion erster Ordnung wir genommen haben. Ist $n$ positiv, so wird $\xi$ Null, ist $n$ negativ, so wird $\xi$ unendlich, und ist $n = 0$, so hat $\xi$ in $\mathfrak{P}$ einen endlichen, von Null verschiedenen Wert.

Die Definition der Ordnungszahl haftet also am Punkte $\mathfrak{P}$. Sie ist für den Körper $\Omega$ invariant.

Daß jede Funktion in jedem Punkte $\mathfrak{P}$ eine bestimmte ganze Zahl $n$ als Ordnungszahl erhält, ergibt sich nun aus \S 185, (5). Man nehme eine Funktion $z$, die in $\mathfrak{P}$ endlich bleibt, als unabhängige Variable, bestimme das zu $\mathfrak{P}$ gehörige Primfunktional $\pi$ und wähle eine Funktion $\omega$, deren Zähler durch $\pi$, aber nicht durch $\pi^2$ teilbar ist. Diese Funktion ist von der ersten Ordnung.

Dann kann man jede Funktion $\xi$ in die Form setzen:
\begin{equation}
\xi = a\omega^n + \omega^{n+1}\eta_1,
\end{equation}
worin $a$ eine von Null verschiedene Konstante ist, und $n$ ist die Ordnungszahl von $\xi$.

Hat $\omega$ die Ordnung $1$, so hat $\omega^n$ die Ordnung $n$.

Über die Ordnungszahlen gelten folgende Sätze, die alle leicht aus (1) folgen:

\begin{satz}
Die Ordnung eines Produktes zweier Funktionen ist gleich der Summe der Ordnungen der Faktoren. Die Ordnung eines Quotienten ist gleich der Differenz der Ordnungen von Zähler und Nenner. Die Ordnung einer Summe ist gleich der niedrigsten unter den Ordnungen der Summanden. Kommen unter den Summanden mehrere von gleicher niedrigster Ordnungszahl vor, so kann die Ordnung der Summe größer, aber nicht kleiner sein.
\end{satz}

Von einer Funktion, die in $\mathfrak{P}$ Null in der $n$ten Ordnung wird, sagt man auch, sie wird Null in der ersten Ordnung in $n$ zusammengefallenen Punkten.

\newpage
\section*{\S 189. Polygone.}
\addcontentsline{toc}{section}{\S 189. Polygone.}

\begin{defin}
Komplexe von Punkten, die den selben Punkt auch mehrmals enthalten können, heißen \emph{Polygone} oder \emph{Vielecke} (des Körpers $\Omega$).
\end{defin}

Die Zahl der Punkte eines Polygons heißt seine Ordnung, ein Polygon $m$ter Ordnung wird auch kurz ein $m$-Eck genannt.

Als Bezeichnung für die Polygone sollen die Buchstaben des großen deutschen Alphabets $\mathfrak{A}$, $\mathfrak{B}$, $\mathfrak{C}$, $\ldots$ dienen.

\begin{defin}
Unter dem Produkt $\mathfrak{A}\mathfrak{B}$ zweier Polygone versteht man das Polygon, das die Punkte von $\mathfrak{A}$ und $\mathfrak{B}$ zugleich enthält, und zwar einen Punkt $\mathfrak{P}$, der $r$mal in $\mathfrak{A}$ und $s$mal in $\mathfrak{B}$ vorkommt, $r + s$mal. Die Ordnung eines Produktes ist gleich der Summe der Ordnungen der Faktoren.
\end{defin}

Wenn wir also mit $\mathfrak{P}_1$, $\mathfrak{P}_2$, $\ldots$ verschiedene Punkte bezeichnen, so können wir jedes Polygon $\mathfrak{A}$ auf eine Weise in die Form setzen:
\begin{equation}
\mathfrak{A} = \mathfrak{P}_1^{r_1} \mathfrak{P}_2^{r_2} \cdots \mathfrak{P}_s^{r_s}.
\end{equation}
Die Punkte sind also in der Rechnung mit Polygonen die Primelemente.

Um auch für die Einheit einen Vertreter zu haben, muß noch das \glqq Nulleck`` $\mathfrak{O}$, das gar keinen Punkt enthält, mitgenannt werden.

Der größte gemeinschaftliche Teiler zweier Polygone $\mathfrak{A}$, $\mathfrak{B}$ enthält jeden Punkt, der in $\mathfrak{A}$ und $\mathfrak{B}$ vorkommt, und zwar so oft, als er in jedem von beiden vorkommt.

Das kleinste gemeinschaftliche Vielfache enthält jeden Punkt von $\mathfrak{A}$ und $\mathfrak{B}$, und zwar in der höchsten der Ordnungen, in denen er in $\mathfrak{A}$ und $\mathfrak{B}$ vorkommt.

\begin{satz}
Ist $z$ eine Variable des Körpers $\Omega$ und $n$ der Grad des Körpers in bezug auf $z$, so nimmt $z$ jeden Wert $c$ in $n$ Punkten an.
\end{satz}

Denn zerlegt man $z - c$ in seine Primfaktoren [\S 183, (17)]:
\[z - c = \pi_1^{e_1} \pi_2^{e_2} \cdots \pi_h^{e_h} \cdot E,\]
so erzeugt jeder dieser Primfaktoren einen Punkt, z.B. $\pi_i$ den Punkt $\mathfrak{P}_i$, in dem die Funktion $z - c$ Null in der $e_i$ten Ordnung wird. Rechnet man diesen Punkt $e_i$fach, so erhält man ein Polygon $\mathfrak{P}_1^{e_1} \mathfrak{P}_2^{e_2} \cdots \mathfrak{P}_h^{e_h}$ von der Ordnung $n$, in dessen Punkten die Funktion $z - c$ verschwindet.

Derselbe Satz gilt aber auch für den Wert $c = \\infty$, wie man daraus schließt, daß der Grad des Körpers $\Omega$ in bezug auf $1:z$ derselbe ist wie in bezug auf $z$.

\begin{defin}
Die $n$ Punkte, in denen $z$ einen Wert $c$ annimmt, heißen \emph{konjugiert nach $z$}, und die Werte $\eta_1$, $\eta_2$, $\ldots$, $\eta_n$, die eine Funktion $\eta$ in $\Omega$ in $n$ konjugierten Punkten annimmt, heißen gleichfalls \emph{konjugiert nach $z$}, und $z$ heißt \emph{von der Ordnung $n$}. Die Konstanten, und nur diese, haben die Ordnung $0$.
\end{defin}

\end{document}
